\documentclass[
  doc,
  floatsintext,
  longtable,
  a4paper,
  nolmodern,
  notxfonts,
  notimes,
  12pt,
  colorlinks=true,linkcolor=blue,citecolor=blue,urlcolor=blue]{apa7}

\usepackage{amsmath}
\usepackage{amssymb}

\geometry{inner=1in, outer=1in}
\fancyhfoffset[LE,RO]{0cm}


\usepackage[bidi=default]{babel}
\babelprovide[main,import]{spanish}


% get rid of language-specific shorthands (see #6817):
\let\LanguageShortHands\languageshorthands
\def\languageshorthands#1{}

\RequirePackage{longtable}
\RequirePackage{threeparttablex}

\makeatletter
\renewcommand{\paragraph}{\@startsection{paragraph}{4}{\parindent}%
	{0\baselineskip \@plus 0.2ex \@minus 0.2ex}%
	{-.5em}%
	{\normalfont\normalsize\bfseries\typesectitle}}

\renewcommand{\subparagraph}[1]{\@startsection{subparagraph}{5}{0.5em}%
	{0\baselineskip \@plus 0.2ex \@minus 0.2ex}%
	{-\z@\relax}%
	{\normalfont\normalsize\bfseries\itshape\hspace{\parindent}{#1}\textit{\addperi}}{\relax}}
\makeatother




\usepackage{longtable, booktabs, multirow, multicol, colortbl, hhline, caption, array, float, xpatch}
\usepackage{subcaption}


\renewcommand\thesubfigure{\Alph{subfigure}}
\setcounter{topnumber}{2}
\setcounter{bottomnumber}{2}
\setcounter{totalnumber}{4}
\renewcommand{\topfraction}{0.85}
\renewcommand{\bottomfraction}{0.85}
\renewcommand{\textfraction}{0.15}
\renewcommand{\floatpagefraction}{0.7}

\usepackage{tcolorbox}
\tcbuselibrary{listings,theorems, breakable, skins}
\usepackage{fontawesome5}

\definecolor{quarto-callout-color}{HTML}{909090}
\definecolor{quarto-callout-note-color}{HTML}{0758E5}
\definecolor{quarto-callout-important-color}{HTML}{CC1914}
\definecolor{quarto-callout-warning-color}{HTML}{EB9113}
\definecolor{quarto-callout-tip-color}{HTML}{00A047}
\definecolor{quarto-callout-caution-color}{HTML}{FC5300}
\definecolor{quarto-callout-color-frame}{HTML}{ACACAC}
\definecolor{quarto-callout-note-color-frame}{HTML}{4582EC}
\definecolor{quarto-callout-important-color-frame}{HTML}{D9534F}
\definecolor{quarto-callout-warning-color-frame}{HTML}{F0AD4E}
\definecolor{quarto-callout-tip-color-frame}{HTML}{02B875}
\definecolor{quarto-callout-caution-color-frame}{HTML}{FD7E14}

%\newlength\Oldarrayrulewidth
%\newlength\Oldtabcolsep


\usepackage{hyperref}



\usepackage{color}
\usepackage{fancyvrb}
\newcommand{\VerbBar}{|}
\newcommand{\VERB}{\Verb[commandchars=\\\{\}]}
\DefineVerbatimEnvironment{Highlighting}{Verbatim}{commandchars=\\\{\}}
% Add ',fontsize=\small' for more characters per line
\usepackage{framed}
\definecolor{shadecolor}{RGB}{241,243,245}
\newenvironment{Shaded}{\begin{snugshade}}{\end{snugshade}}
\newcommand{\AlertTok}[1]{\textcolor[rgb]{0.68,0.00,0.00}{#1}}
\newcommand{\AnnotationTok}[1]{\textcolor[rgb]{0.37,0.37,0.37}{#1}}
\newcommand{\AttributeTok}[1]{\textcolor[rgb]{0.40,0.45,0.13}{#1}}
\newcommand{\BaseNTok}[1]{\textcolor[rgb]{0.68,0.00,0.00}{#1}}
\newcommand{\BuiltInTok}[1]{\textcolor[rgb]{0.00,0.23,0.31}{#1}}
\newcommand{\CharTok}[1]{\textcolor[rgb]{0.13,0.47,0.30}{#1}}
\newcommand{\CommentTok}[1]{\textcolor[rgb]{0.37,0.37,0.37}{#1}}
\newcommand{\CommentVarTok}[1]{\textcolor[rgb]{0.37,0.37,0.37}{\textit{#1}}}
\newcommand{\ConstantTok}[1]{\textcolor[rgb]{0.56,0.35,0.01}{#1}}
\newcommand{\ControlFlowTok}[1]{\textcolor[rgb]{0.00,0.23,0.31}{\textbf{#1}}}
\newcommand{\DataTypeTok}[1]{\textcolor[rgb]{0.68,0.00,0.00}{#1}}
\newcommand{\DecValTok}[1]{\textcolor[rgb]{0.68,0.00,0.00}{#1}}
\newcommand{\DocumentationTok}[1]{\textcolor[rgb]{0.37,0.37,0.37}{\textit{#1}}}
\newcommand{\ErrorTok}[1]{\textcolor[rgb]{0.68,0.00,0.00}{#1}}
\newcommand{\ExtensionTok}[1]{\textcolor[rgb]{0.00,0.23,0.31}{#1}}
\newcommand{\FloatTok}[1]{\textcolor[rgb]{0.68,0.00,0.00}{#1}}
\newcommand{\FunctionTok}[1]{\textcolor[rgb]{0.28,0.35,0.67}{#1}}
\newcommand{\ImportTok}[1]{\textcolor[rgb]{0.00,0.46,0.62}{#1}}
\newcommand{\InformationTok}[1]{\textcolor[rgb]{0.37,0.37,0.37}{#1}}
\newcommand{\KeywordTok}[1]{\textcolor[rgb]{0.00,0.23,0.31}{\textbf{#1}}}
\newcommand{\NormalTok}[1]{\textcolor[rgb]{0.00,0.23,0.31}{#1}}
\newcommand{\OperatorTok}[1]{\textcolor[rgb]{0.37,0.37,0.37}{#1}}
\newcommand{\OtherTok}[1]{\textcolor[rgb]{0.00,0.23,0.31}{#1}}
\newcommand{\PreprocessorTok}[1]{\textcolor[rgb]{0.68,0.00,0.00}{#1}}
\newcommand{\RegionMarkerTok}[1]{\textcolor[rgb]{0.00,0.23,0.31}{#1}}
\newcommand{\SpecialCharTok}[1]{\textcolor[rgb]{0.37,0.37,0.37}{#1}}
\newcommand{\SpecialStringTok}[1]{\textcolor[rgb]{0.13,0.47,0.30}{#1}}
\newcommand{\StringTok}[1]{\textcolor[rgb]{0.13,0.47,0.30}{#1}}
\newcommand{\VariableTok}[1]{\textcolor[rgb]{0.07,0.07,0.07}{#1}}
\newcommand{\VerbatimStringTok}[1]{\textcolor[rgb]{0.13,0.47,0.30}{#1}}
\newcommand{\WarningTok}[1]{\textcolor[rgb]{0.37,0.37,0.37}{\textit{#1}}}

\providecommand{\tightlist}{%
  \setlength{\itemsep}{0pt}\setlength{\parskip}{0pt}}
\usepackage{longtable,booktabs,array}
\newcounter{none} % for unnumbered tables
\usepackage{calc} % for calculating minipage widths
% Correct order of tables after \paragraph or \subparagraph
\usepackage{etoolbox}
\makeatletter
\patchcmd\longtable{\par}{\if@noskipsec\mbox{}\fi\par}{}{}
\makeatother
% Allow footnotes in longtable head/foot
\IfFileExists{footnotehyper.sty}{\usepackage{footnotehyper}}{\usepackage{footnote}}
\makesavenoteenv{longtable}

\usepackage{graphicx}
\makeatletter
\newsavebox\pandoc@box
\newcommand*\pandocbounded[1]{% scales image to fit in text height/width
  \sbox\pandoc@box{#1}%
  \Gscale@div\@tempa{\textheight}{\dimexpr\ht\pandoc@box+\dp\pandoc@box\relax}%
  \Gscale@div\@tempb{\linewidth}{\wd\pandoc@box}%
  \ifdim\@tempb\p@<\@tempa\p@\let\@tempa\@tempb\fi% select the smaller of both
  \ifdim\@tempa\p@<\p@\scalebox{\@tempa}{\usebox\pandoc@box}%
  \else\usebox{\pandoc@box}%
  \fi%
}
% Set default figure placement to htbp
\def\fps@figure{htbp}
\makeatother







\usepackage{newtx}

\defaultfontfeatures{Scale=MatchLowercase}
\defaultfontfeatures[\rmfamily]{Ligatures=TeX,Scale=1}





\title{Scrcpy}


\shorttitle{Scrcpy}


\usepackage{etoolbox}



\ccoppy{\textcopyright~2021}



\author{Edison Achalma}



\affiliation{
{Escuela Profesional de Economía, Universidad Nacional de San Cristóbal
de Huamanga}}




\leftheader{Achalma}

\date{2021-10-21}


\abstract{Este abstract será actualizado una vez que se complete el
contenido final del artículo. }

\keywords{keyword1, keyword2}

\authornote{\par{\addORCIDlink{Edison Achalma}{0000-0001-6996-3364}} 
\par{ }
\par{   El autor no tiene conflictos de interés que revelar.    Los
roles de autor se clasificaron utilizando la taxonomía de roles de
colaborador (CRediT; https://credit.niso.org/) de la siguiente
manera:  Edison Achalma:   conceptualización, metodología, análisis
formal, investigación, recursos, curación de
datos, redacción, visualización, supervisión, administración del
proyecto}
\par{La correspondencia relativa a este artículo debe dirigirse a Edison
Achalma, Escuela Profesional de Economía, Universidad Nacional de San
Cristóbal de Huamanga, Portal Independencia N°
57, Ayacucho, AYA 5001, Perú, Email: \href{mailto:elmer.achalma.09@unsch.edu.pe}{elmer.achalma.09@unsch.edu.pe}}
}

\makeatletter
\let\endoldlt\endlongtable
\def\endlongtable{
\hline
\endoldlt
}
\makeatother

\urlstyle{same}



\hypersetup{
  pdftitle={Scrcpy},
  pdfauthor={Autor}
}
\makeatletter
\@ifpackageloaded{caption}{}{\usepackage{caption}}
\AtBeginDocument{%
\ifdefined\contentsname
  \renewcommand*\contentsname{Tabla de contenidos}
\else
  \newcommand\contentsname{Tabla de contenidos}
\fi
\ifdefined\listfigurename
  \renewcommand*\listfigurename{Lista de Figuras}
\else
  \newcommand\listfigurename{Lista de Figuras}
\fi
\ifdefined\listtablename
  \renewcommand*\listtablename{Lista de Tablas}
\else
  \newcommand\listtablename{Lista de Tablas}
\fi
\ifdefined\figurename
  \renewcommand*\figurename{Figura}
\else
  \newcommand\figurename{Figura}
\fi
\ifdefined\tablename
  \renewcommand*\tablename{Tabla}
\else
  \newcommand\tablename{Tabla}
\fi
}
\@ifpackageloaded{float}{}{\usepackage{float}}
\floatstyle{ruled}
\@ifundefined{c@chapter}{\newfloat{codelisting}{h}{lop}}{\newfloat{codelisting}{h}{lop}[chapter]}
\floatname{codelisting}{Listado}
\newcommand*\listoflistings{\listof{codelisting}{Lista de Listados}}
\makeatother
\makeatletter
\makeatother
\makeatletter
\@ifpackageloaded{caption}{}{\usepackage{caption}}
\@ifpackageloaded{subcaption}{}{\usepackage{subcaption}}
\makeatother
\makeatletter
\@ifpackageloaded{fontawesome5}{}{\usepackage{fontawesome5}}
\makeatother

% From https://tex.stackexchange.com/a/645996/211326
%%% apa7 doesn't want to add appendix section titles in the toc
%%% let's make it do it
\makeatletter
\xpatchcmd{\appendix}
  {\par}
  {\addcontentsline{toc}{section}{\@currentlabelname}\par}
  {}{}
\makeatother

%% Disable longtable counter
%% https://tex.stackexchange.com/a/248395/211326

\usepackage{etoolbox}

\makeatletter
\patchcmd{\LT@caption}
  {\bgroup}
  {\bgroup\global\LTpatch@captiontrue}
  {}{}
\patchcmd{\longtable}
  {\par}
  {\par\global\LTpatch@captionfalse}
  {}{}
\apptocmd{\endlongtable}
  {\ifLTpatch@caption\else\addtocounter{table}{-1}\fi}
  {}{}
\newif\ifLTpatch@caption
\makeatother

\begin{document}

\maketitle


\hypertarget{toc}{}
\tableofcontents
\newpage
\section[Introduction]{Scrcpy}

\setcounter{secnumdepth}{3}

\setlength\LTleft{0pt}




\section{Scrcpy}\label{scrcpy}

\subsection{¿Qué es scrcpy?}\label{quuxe9-es-scrcpy}

\textbf{scrcpy} (screen copy) es una aplicación de código abierto que
permite:

\begin{itemize}
\tightlist
\item
  \textbf{Duplicar} (mirror) la pantalla de dispositivos Android en tu
  computadora
\item
  \textbf{Controlar} el dispositivo desde el teclado y mouse de la
  computadora
\item
  \textbf{Grabar} video y audio del dispositivo
\item
  \textbf{Transmitir} audio desde el dispositivo
\end{itemize}

\textbf{Pronunciación:} ``screen copy''

\subsection{Filosofía de Diseño}\label{filosofuxeda-de-diseuxf1o}

scrcpy se enfoca en:

\begin{enumerate}
\def\labelenumi{\arabic{enumi}.}
\tightlist
\item
  \textbf{Ligereza:} Aplicación nativa, solo muestra la pantalla del
  dispositivo
\item
  \textbf{Rendimiento:} 30-120 fps dependiendo del dispositivo
\item
  \textbf{Calidad:} Hasta 1920×1080 o superior
\item
  \textbf{Baja latencia:} 35-70ms
\item
  \textbf{Inicio rápido:} \textasciitilde1 segundo para mostrar la
  primera imagen
\item
  \textbf{No intrusivo:} No instala nada en el dispositivo Android
\item
  \textbf{Sin publicidad:} Software libre y de código abierto
\end{enumerate}

\subsection{Ventajas sobre Otras
Soluciones}\label{ventajas-sobre-otras-soluciones}

{\def\LTcaptype{none} % do not increment counter
\begin{longtable}[]{@{}lllll@{}}
\toprule\noalign{}
Característica & scrcpy & Vysor & AirDroid & TeamViewer \\
\midrule\noalign{}
\endhead
\bottomrule\noalign{}
\endlastfoot
Gratuito & Sí & Limitado & Limitado & No \\
Código abierto & Sí & No & No & No \\
Baja latencia & Sí & No & No & No \\
Sin anuncios & Sí & No & No & Sí \\
Sin cuenta & Sí & No & No & No \\
Funciona sin internet & Sí & Sí & No & No \\
Grabación nativa & Sí & Premium & Premium & No \\
\end{longtable}
}

\section{Características
Principales}\label{caracteruxedsticas-principales}

\subsection{Funcionalidades
Disponibles}\label{funcionalidades-disponibles}

\begin{itemize}
\tightlist
\item
  \textbf{Duplicación de pantalla} (mirroring) en tiempo real
\item
  \textbf{Reenvío de audio} (Android 11+)
\item
  \textbf{Grabación} de video y audio
\item
  \textbf{Pantalla virtual} (virtual display)
\item
  \textbf{Duplicación con pantalla apagada} del dispositivo
\item
  \textbf{Copiar-pegar} bidireccional (texto e imágenes)
\item
  \textbf{Calidad configurable} de video
\item
  \textbf{Duplicación de cámara} (Android 12+)
\item
  \textbf{Webcam virtual} (V4L2 en Linux)
\item
  \textbf{Simulación de teclado y mouse físicos} (HID)
\item
  \textbf{Soporte de gamepad}
\item
  \textbf{Modo OTG} (control sin depuración USB)
\end{itemize}

\subsection{Casos de Uso Ideales}\label{casos-de-uso-ideales}

\textbf{Para ti como economista e informático:}

\begin{enumerate}
\def\labelenumi{\arabic{enumi}.}
\item
  \textbf{Presentaciones Académicas}

  \begin{itemize}
  \tightlist
  \item
    Mostrar aplicaciones de econometría móviles
  \item
    Demostrar visualizaciones de datos
  \item
    Presentar simulaciones interactivas
  \end{itemize}
\item
  \textbf{Tutoriales y Documentación}

  \begin{itemize}
  \tightlist
  \item
    Grabar tutoriales de aplicaciones
  \item
    Crear material educativo
  \item
    Documentar procesos de análisis
  \end{itemize}
\item
  \textbf{Desarrollo y Testing}

  \begin{itemize}
  \tightlist
  \item
    Probar aplicaciones econométricas
  \item
    Verificar visualizaciones en móvil
  \item
    Debug de aplicaciones web responsive
  \end{itemize}
\item
  \textbf{Colaboración Remota}

  \begin{itemize}
  \tightlist
  \item
    Compartir pantalla en reuniones
  \item
    Asistencia técnica remota
  \item
    Demostraciones en vivo
  \end{itemize}
\end{enumerate}

\section{Requisitos Previos}\label{requisitos-previos}

\subsection{En el Dispositivo Android}\label{en-el-dispositivo-android}

\textbf{Requisitos mínimos:}

\begin{itemize}
\tightlist
\item
  Android 5.0 (API 21) o superior
\item
  Para audio: Android 11 (API 30) o superior
\end{itemize}

\textbf{Configuración necesaria:}

\begin{enumerate}
\def\labelenumi{\arabic{enumi}.}
\item
  \textbf{Habilitar Depuración USB}

  Ir a: \texttt{Ajustes} → \texttt{Acerca\ del\ teléfono} → Tocar 7
  veces en \texttt{Número\ de\ compilación}

  Luego: \texttt{Ajustes} → \texttt{Opciones\ de\ desarrollador} →
  Activar \texttt{Depuración\ USB}
\item
  \textbf{Dispositivos Xiaomi (adicional)}

  Algunos dispositivos Xiaomi requieren:

  \begin{itemize}
  \tightlist
  \item
    \texttt{Depuración\ USB\ (Configuración\ de\ seguridad)} activado
  \item
    Reiniciar el dispositivo después de activarlo
  \end{itemize}
\end{enumerate}

\textbf{Verificar que funciona:}

\begin{Shaded}
\begin{Highlighting}[]
\CommentTok{\# Conectar dispositivo por USB}
\CommentTok{\# Verificar que aparece}
\ExtensionTok{adb}\NormalTok{ devices}
\end{Highlighting}
\end{Shaded}

\subsection{En Arch Linux}\label{en-arch-linux}

\textbf{Paquetes necesarios:}

\begin{Shaded}
\begin{Highlighting}[]
\CommentTok{\# Verificar que tienes instalado}
\ExtensionTok{pacman} \AttributeTok{{-}Q}\NormalTok{ ffmpeg sdl2 adb libusb}
\end{Highlighting}
\end{Shaded}

\textbf{Si falta algo:}

\begin{Shaded}
\begin{Highlighting}[]
\FunctionTok{sudo}\NormalTok{ pacman }\AttributeTok{{-}S}\NormalTok{ ffmpeg sdl2 android{-}tools libusb}
\end{Highlighting}
\end{Shaded}

\begin{center}\rule{0.5\linewidth}{0.5pt}\end{center}

\section{Instalación en Arch Linux}\label{instalaciuxf3n-en-arch-linux}

\subsection{Método 1: Instalación desde Repositorios Oficiales
(Recomendado)}\label{muxe9todo-1-instalaciuxf3n-desde-repositorios-oficiales-recomendado}

\begin{Shaded}
\begin{Highlighting}[]
\CommentTok{\# scrcpy está en los repositorios oficiales de Arch}
\FunctionTok{sudo}\NormalTok{ pacman }\AttributeTok{{-}S}\NormalTok{ scrcpy}
\end{Highlighting}
\end{Shaded}

\textbf{Verificar instalación:}

\begin{Shaded}
\begin{Highlighting}[]
\ExtensionTok{scrcpy} \AttributeTok{{-}{-}version}
\CommentTok{\# Debe mostrar: scrcpy 3.3.4 o superior}
\end{Highlighting}
\end{Shaded}

\subsection{Método 2: Instalación desde AUR (Versión de
Desarrollo)}\label{muxe9todo-2-instalaciuxf3n-desde-aur-versiuxf3n-de-desarrollo}

\begin{Shaded}
\begin{Highlighting}[]
\CommentTok{\# Si quieres la versión más reciente (desarrollo)}
\ExtensionTok{yay} \AttributeTok{{-}S}\NormalTok{ scrcpy{-}git}

\CommentTok{\# O con paru}
\ExtensionTok{paru} \AttributeTok{{-}S}\NormalTok{ scrcpy{-}git}
\end{Highlighting}
\end{Shaded}

\subsection{Método 3: Compilación desde Código
Fuente}\label{muxe9todo-3-compilaciuxf3n-desde-cuxf3digo-fuente}

\textbf{Solo si necesitas modificar el código o tienes requisitos
específicos.}

\begin{Shaded}
\begin{Highlighting}[]
\CommentTok{\# 1. Instalar dependencias de compilación}
\FunctionTok{sudo}\NormalTok{ pacman }\AttributeTok{{-}S}\NormalTok{ base{-}devel meson ninja pkg{-}config }\DataTypeTok{\textbackslash{}}
\NormalTok{  ffmpeg libusb sdl2 }\DataTypeTok{\textbackslash{}}
\NormalTok{  jdk17{-}openjdk}

\CommentTok{\# 2. Clonar repositorio}
\FunctionTok{git}\NormalTok{ clone https://github.com/Genymobile/scrcpy.git}
\BuiltInTok{cd}\NormalTok{ scrcpy}

\CommentTok{\# 3. Compilar}
\ExtensionTok{meson}\NormalTok{ setup x }\AttributeTok{{-}{-}buildtype}\OperatorTok{=}\NormalTok{release }\AttributeTok{{-}{-}strip} \AttributeTok{{-}Db\_lto}\OperatorTok{=}\NormalTok{true}
\ExtensionTok{ninja} \AttributeTok{{-}Cx}

\CommentTok{\# 4. Instalar}
\FunctionTok{sudo}\NormalTok{ ninja }\AttributeTok{{-}Cx}\NormalTok{ install}

\CommentTok{\# 5. Verificar}
\ExtensionTok{scrcpy} \AttributeTok{{-}{-}version}
\end{Highlighting}
\end{Shaded}

\section{Configuración Inicial}\label{configuraciuxf3n-inicial}

\subsection{1. Configurar ADB}\label{configurar-adb}

\textbf{ADB (Android Debug Bridge) es necesario para scrcpy.}

\begin{Shaded}
\begin{Highlighting}[]
\CommentTok{\# Verificar que adb funciona}
\ExtensionTok{adb} \AttributeTok{{-}{-}version}

\CommentTok{\# Iniciar servidor adb}
\ExtensionTok{adb}\NormalTok{ start{-}server}

\CommentTok{\# Listar dispositivos conectados}
\ExtensionTok{adb}\NormalTok{ devices}

\CommentTok{\# Debe mostrar algo como:}
\CommentTok{\# List of devices attached}
\CommentTok{\# 0123456789ABCDEF    device}
\end{Highlighting}
\end{Shaded}

\subsection{2. Configurar Permisos de
Usuario}\label{configurar-permisos-de-usuario}

\textbf{En Arch Linux, tu usuario debe estar en el grupo
\texttt{adbusers}:}

\begin{Shaded}
\begin{Highlighting}[]
\CommentTok{\# Verificar si estás en el grupo}
\FunctionTok{groups} \KeywordTok{|} \FunctionTok{grep}\NormalTok{ adbusers}

\CommentTok{\# Si no estás, agregarte}
\FunctionTok{sudo}\NormalTok{ usermod }\AttributeTok{{-}aG}\NormalTok{ adbusers }\VariableTok{$USER}

\CommentTok{\# Cerrar sesión y volver a iniciar}
\CommentTok{\# O ejecutar:}
\ExtensionTok{newgrp}\NormalTok{ adbusers}
\end{Highlighting}
\end{Shaded}

\subsection{3. Configurar Reglas udev}\label{configurar-reglas-udev}

\textbf{Para que adb detecte dispositivos sin sudo:}

\begin{Shaded}
\begin{Highlighting}[]
\CommentTok{\# Crear archivo de reglas}
\FunctionTok{sudo}\NormalTok{ nano /etc/udev/rules.d/51{-}android.rules}
\end{Highlighting}
\end{Shaded}

\textbf{Agregar (ejemplo para varios fabricantes):}

\begin{verbatim}
# Google
SUBSYSTEM=="usb", ATTR{idVendor}=="18d1", MODE="0666", GROUP="adbusers"
# Samsung
SUBSYSTEM=="usb", ATTR{idVendor}=="04e8", MODE="0666", GROUP="adbusers"
# Xiaomi
SUBSYSTEM=="usb", ATTR{idVendor}=="2717", MODE="0666", GROUP="adbusers"
# Huawei
SUBSYSTEM=="usb", ATTR{idVendor}=="12d1", MODE="0666", GROUP="adbusers"
# OnePlus
SUBSYSTEM=="usb", ATTR{idVendor}=="2a70", MODE="0666", GROUP="adbusers"
\end{verbatim}

\textbf{Recargar reglas:}

\begin{Shaded}
\begin{Highlighting}[]
\FunctionTok{sudo}\NormalTok{ udevadm control }\AttributeTok{{-}{-}reload{-}rules}
\FunctionTok{sudo}\NormalTok{ udevadm trigger}
\end{Highlighting}
\end{Shaded}

\subsection{4. Primera Conexión}\label{primera-conexiuxf3n}

\begin{Shaded}
\begin{Highlighting}[]
\CommentTok{\# Conectar dispositivo por USB}
\CommentTok{\# En el dispositivo, autorizar la computadora (popup que aparece)}

\CommentTok{\# Verificar conexión}
\ExtensionTok{adb}\NormalTok{ devices}

\CommentTok{\# Iniciar scrcpy}
\ExtensionTok{scrcpy}
\end{Highlighting}
\end{Shaded}

\textbf{Debe aparecer la pantalla del dispositivo en una ventana.}

\section{Uso Básico}\label{uso-buxe1sico}

\subsection{Comando Básico}\label{comando-buxe1sico}

\begin{Shaded}
\begin{Highlighting}[]
\CommentTok{\# Duplicar pantalla del dispositivo}
\ExtensionTok{scrcpy}
\end{Highlighting}
\end{Shaded}

\textbf{Eso es todo.} scrcpy automáticamente:

\begin{enumerate}
\def\labelenumi{\arabic{enumi}.}
\tightlist
\item
  Detecta el dispositivo conectado
\item
  Instala el servidor temporal en el dispositivo
\item
  Inicia la duplicación de pantalla
\item
  Habilita el control con teclado y mouse
\end{enumerate}

\subsection{Opciones Más Comunes}\label{opciones-muxe1s-comunes}

\begin{Shaded}
\begin{Highlighting}[]
\CommentTok{\# Limitar resolución (mejor rendimiento)}
\ExtensionTok{scrcpy} \AttributeTok{{-}m1920}        \CommentTok{\# máximo 1920px de ancho}
\ExtensionTok{scrcpy} \AttributeTok{{-}m1024}        \CommentTok{\# máximo 1024px (más fluido)}

\CommentTok{\# Solo video, sin audio}
\ExtensionTok{scrcpy} \AttributeTok{{-}{-}no{-}audio}

\CommentTok{\# Solo audio, sin video}
\ExtensionTok{scrcpy} \AttributeTok{{-}{-}no{-}video} \AttributeTok{{-}{-}no{-}control}

\CommentTok{\# Ventana siempre visible}
\ExtensionTok{scrcpy} \AttributeTok{{-}{-}always{-}on{-}top}

\CommentTok{\# Pantalla completa}
\ExtensionTok{scrcpy} \AttributeTok{{-}f}

\CommentTok{\# Mantener dispositivo despierto (conectado a corriente)}
\ExtensionTok{scrcpy} \AttributeTok{{-}w}

\CommentTok{\# Apagar pantalla del dispositivo (ahorra batería)}
\ExtensionTok{scrcpy} \AttributeTok{{-}S}
\end{Highlighting}
\end{Shaded}

\subsection{Ejemplos de Uso Rápido}\label{ejemplos-de-uso-ruxe1pido}

\begin{Shaded}
\begin{Highlighting}[]
\CommentTok{\# Presentación fluida (baja resolución, alta velocidad)}
\ExtensionTok{scrcpy} \AttributeTok{{-}m1024} \AttributeTok{{-}{-}max{-}fps}\OperatorTok{=}\NormalTok{60 }\AttributeTok{{-}{-}no{-}audio}

\CommentTok{\# Grabación de alta calidad}
\ExtensionTok{scrcpy} \AttributeTok{{-}{-}video{-}codec}\OperatorTok{=}\NormalTok{h265 }\AttributeTok{{-}m1920} \AttributeTok{{-}{-}max{-}fps}\OperatorTok{=}\NormalTok{60 }\AttributeTok{{-}r}\NormalTok{ tutorial.mp4}

\CommentTok{\# Solo audio (como dictáfono)}
\ExtensionTok{scrcpy} \AttributeTok{{-}{-}audio{-}source}\OperatorTok{=}\NormalTok{mic }\AttributeTok{{-}{-}no{-}video} \AttributeTok{{-}{-}no{-}control} \AttributeTok{{-}r}\NormalTok{ grabacion.opus}

\CommentTok{\# Control sin duplicación (ahorra recursos)}
\ExtensionTok{scrcpy} \AttributeTok{{-}{-}no{-}video} \AttributeTok{{-}{-}no{-}audio}
\end{Highlighting}
\end{Shaded}

\section{Conexión del Dispositivo}\label{conexiuxf3n-del-dispositivo}

\subsection{Selección de Dispositivo}\label{selecciuxf3n-de-dispositivo}

\textbf{Si solo hay un dispositivo conectado:}

\begin{Shaded}
\begin{Highlighting}[]
\ExtensionTok{scrcpy}  \CommentTok{\# Se conecta automáticamente}
\end{Highlighting}
\end{Shaded}

\textbf{Si hay múltiples dispositivos:}

\begin{Shaded}
\begin{Highlighting}[]
\CommentTok{\# Listar dispositivos}
\ExtensionTok{adb}\NormalTok{ devices}

\CommentTok{\# Conectar por serial}
\ExtensionTok{scrcpy} \AttributeTok{{-}{-}serial}\OperatorTok{=}\NormalTok{0123456789ABCDEF}
\ExtensionTok{scrcpy} \AttributeTok{{-}s}\NormalTok{ 0123456789ABCDEF  }\CommentTok{\# versión corta}

\CommentTok{\# Conectar al único dispositivo USB}
\ExtensionTok{scrcpy} \AttributeTok{{-}{-}select{-}usb}
\ExtensionTok{scrcpy} \AttributeTok{{-}d}  \CommentTok{\# versión corta}

\CommentTok{\# Conectar al único dispositivo TCP/IP}
\ExtensionTok{scrcpy} \AttributeTok{{-}{-}select{-}tcpip}
\ExtensionTok{scrcpy} \AttributeTok{{-}e}  \CommentTok{\# versión corta}
\end{Highlighting}
\end{Shaded}

\subsection{Conexión TCP/IP
(Inalámbrica)}\label{conexiuxf3n-tcpip-inaluxe1mbrica}

\subsubsection{Método 1: Automático}\label{muxe9todo-1-automuxe1tico}

\textbf{Desde USB a WiFi automáticamente:}

\begin{Shaded}
\begin{Highlighting}[]
\CommentTok{\# Conectar dispositivo por USB primero}
\ExtensionTok{scrcpy} \AttributeTok{{-}{-}tcpip}

\CommentTok{\# scrcpy detecta la IP, habilita TCP/IP, y se conecta}
\CommentTok{\# Ahora puedes desconectar el cable USB}
\end{Highlighting}
\end{Shaded}

\textbf{Si ya sabes la IP del dispositivo:}

\begin{Shaded}
\begin{Highlighting}[]
\ExtensionTok{scrcpy} \AttributeTok{{-}{-}tcpip}\OperatorTok{=}\NormalTok{192.168.1.100       }\CommentTok{\# puerto por defecto 5555}
\ExtensionTok{scrcpy} \AttributeTok{{-}{-}tcpip}\OperatorTok{=}\NormalTok{192.168.1.100:5555  }\CommentTok{\# especificar puerto}
\end{Highlighting}
\end{Shaded}

\subsubsection{Método 2: Manual}\label{muxe9todo-2-manual}

\begin{Shaded}
\begin{Highlighting}[]
\CommentTok{\# 1. Conectar dispositivo por USB}
\CommentTok{\# 2. Habilitar TCP/IP en el dispositivo}
\ExtensionTok{adb}\NormalTok{ tcpip 5555}

\CommentTok{\# 3. Obtener IP del dispositivo}
\ExtensionTok{adb}\NormalTok{ shell ip route }\KeywordTok{|} \FunctionTok{awk} \StringTok{\textquotesingle{}\{print $9\}\textquotesingle{}}
\CommentTok{\# O desde el dispositivo: Ajustes → Acerca del teléfono → Estado}

\CommentTok{\# 4. Desconectar USB}
\CommentTok{\# 5. Conectar por WiFi}
\ExtensionTok{adb}\NormalTok{ connect 192.168.1.100:5555}

\CommentTok{\# 6. Usar scrcpy normalmente}
\ExtensionTok{scrcpy}

\CommentTok{\# 7. Desconectar cuando termines}
\ExtensionTok{adb}\NormalTok{ disconnect}
\end{Highlighting}
\end{Shaded}

\textbf{Nota:} Desde Android 11 hay una opción de depuración inalámbrica
en Opciones de Desarrollador.

\subsection{Variable de Entorno}\label{variable-de-entorno}

\begin{Shaded}
\begin{Highlighting}[]
\CommentTok{\# En Fish}
\BuiltInTok{set} \AttributeTok{{-}gx}\NormalTok{ ANDROID\_SERIAL 0123456789ABCDEF}
\ExtensionTok{scrcpy}

\CommentTok{\# En Bash/Zsh}
\BuiltInTok{export} \VariableTok{ANDROID\_SERIAL}\OperatorTok{=}\NormalTok{0123456789ABCDEF}
\ExtensionTok{scrcpy}
\end{Highlighting}
\end{Shaded}

\section{Control de Video}\label{control-de-video}

\subsection{Resolución y Calidad}\label{resoluciuxf3n-y-calidad}

\subsubsection{Limitar Resolución}\label{limitar-resoluciuxf3n}

\begin{Shaded}
\begin{Highlighting}[]
\CommentTok{\# Limitar ancho/alto máximo}
\ExtensionTok{scrcpy} \AttributeTok{{-}{-}max{-}size}\OperatorTok{=}\NormalTok{1920}
\ExtensionTok{scrcpy} \AttributeTok{{-}m1920}  \CommentTok{\# versión corta}

\CommentTok{\# Ejemplos por caso de uso:}
\ExtensionTok{scrcpy} \AttributeTok{{-}m1024}  \CommentTok{\# Presentaciones fluidas}
\ExtensionTok{scrcpy} \AttributeTok{{-}m1920}  \CommentTok{\# Balance calidad/rendimiento}
\ExtensionTok{scrcpy} \AttributeTok{{-}m2560}  \CommentTok{\# Alta calidad (si el dispositivo lo soporta)}
\end{Highlighting}
\end{Shaded}

\textbf{Regla:} La otra dimensión se calcula automáticamente manteniendo
la relación de aspecto.

\subsubsection{Tasa de Bits}\label{tasa-de-bits}

\begin{Shaded}
\begin{Highlighting}[]
\CommentTok{\# Cambiar bitrate de video (por defecto 8 Mbps)}
\ExtensionTok{scrcpy} \AttributeTok{{-}{-}video{-}bit{-}rate}\OperatorTok{=}\NormalTok{2M}
\ExtensionTok{scrcpy} \AttributeTok{{-}b}\NormalTok{ 2M  }\CommentTok{\# versión corta}

\CommentTok{\# Ejemplos:}
\ExtensionTok{scrcpy} \AttributeTok{{-}b}\NormalTok{ 2M   }\CommentTok{\# Menor calidad, más fluido}
\ExtensionTok{scrcpy} \AttributeTok{{-}b}\NormalTok{ 4M   }\CommentTok{\# Balance}
\ExtensionTok{scrcpy} \AttributeTok{{-}b}\NormalTok{ 8M   }\CommentTok{\# Por defecto}
\ExtensionTok{scrcpy} \AttributeTok{{-}b}\NormalTok{ 12M  }\CommentTok{\# Alta calidad}
\end{Highlighting}
\end{Shaded}

\subsubsection{Tasa de Fotogramas}\label{tasa-de-fotogramas}

\begin{Shaded}
\begin{Highlighting}[]
\CommentTok{\# Limitar FPS}
\ExtensionTok{scrcpy} \AttributeTok{{-}{-}max{-}fps}\OperatorTok{=}\NormalTok{30   }\CommentTok{\# Ahorra batería}
\ExtensionTok{scrcpy} \AttributeTok{{-}{-}max{-}fps}\OperatorTok{=}\NormalTok{60   }\CommentTok{\# Fluido}
\ExtensionTok{scrcpy} \AttributeTok{{-}{-}max{-}fps}\OperatorTok{=}\NormalTok{120  }\CommentTok{\# Muy fluido (si dispositivo lo soporta)}

\CommentTok{\# Mostrar FPS en consola}
\ExtensionTok{scrcpy} \AttributeTok{{-}{-}print{-}fps}

\CommentTok{\# Mostrar/ocultar FPS en tiempo real}
\CommentTok{\# Presionar MOD+i (Alt+i o Super+i)}
\end{Highlighting}
\end{Shaded}

\subsection{Códec de Video}\label{cuxf3dec-de-video}

\begin{Shaded}
\begin{Highlighting}[]
\CommentTok{\# Seleccionar códec}
\ExtensionTok{scrcpy} \AttributeTok{{-}{-}video{-}codec}\OperatorTok{=}\NormalTok{h264  }\CommentTok{\# Por defecto, menor latencia}
\ExtensionTok{scrcpy} \AttributeTok{{-}{-}video{-}codec}\OperatorTok{=}\NormalTok{h265  }\CommentTok{\# Mejor calidad}
\ExtensionTok{scrcpy} \AttributeTok{{-}{-}video{-}codec}\OperatorTok{=}\NormalTok{av1   }\CommentTok{\# Muy nueva, pocos dispositivos}

\CommentTok{\# Listar codificadores disponibles}
\ExtensionTok{scrcpy} \AttributeTok{{-}{-}list{-}encoders}
\end{Highlighting}
\end{Shaded}

\textbf{Recomendaciones:}

\begin{itemize}
\tightlist
\item
  \textbf{H.264:} Mejor latencia, compatible con todos los dispositivos
\item
  \textbf{H.265:} Mejor calidad, mayor compresión
\item
  \textbf{AV1:} Mejor compresión, pero pocos dispositivos lo soportan
\end{itemize}

\subsection{Orientación de Pantalla}\label{orientaciuxf3n-de-pantalla}

\begin{Shaded}
\begin{Highlighting}[]
\CommentTok{\# Orientación de captura (en el dispositivo)}
\ExtensionTok{scrcpy} \AttributeTok{{-}{-}capture{-}orientation}\OperatorTok{=}\NormalTok{0      }\CommentTok{\# Sin rotación}
\ExtensionTok{scrcpy} \AttributeTok{{-}{-}capture{-}orientation}\OperatorTok{=}\NormalTok{90     }\CommentTok{\# 90° horario}
\ExtensionTok{scrcpy} \AttributeTok{{-}{-}capture{-}orientation}\OperatorTok{=}\NormalTok{180    }\CommentTok{\# 180°}
\ExtensionTok{scrcpy} \AttributeTok{{-}{-}capture{-}orientation}\OperatorTok{=}\NormalTok{270    }\CommentTok{\# 270° horario}

\CommentTok{\# Bloquear orientación (no rota aunque gires el dispositivo)}
\ExtensionTok{scrcpy} \AttributeTok{{-}{-}capture{-}orientation}\OperatorTok{=}\NormalTok{@90    }\CommentTok{\# Bloqueada a 90°}

\CommentTok{\# Orientación de visualización (en la computadora)}
\ExtensionTok{scrcpy} \AttributeTok{{-}{-}orientation}\OperatorTok{=}\NormalTok{90             }\CommentTok{\# Rotar visualización 90°}

\CommentTok{\# Rotar contenido por ángulo personalizado}
\ExtensionTok{scrcpy} \AttributeTok{{-}{-}angle}\OperatorTok{=}\NormalTok{23                   }\CommentTok{\# Rotar 23° horario}
\end{Highlighting}
\end{Shaded}

\subsection{Recorte de Pantalla}\label{recorte-de-pantalla}

\begin{Shaded}
\begin{Highlighting}[]
\CommentTok{\# Recortar región específica}
\ExtensionTok{scrcpy} \AttributeTok{{-}{-}crop}\OperatorTok{=}\NormalTok{1224:1440:0:0   }\CommentTok{\# Ancho:Alto:OffsetX:OffsetY}

\CommentTok{\# Ejemplo: recortar solo la mitad superior}
\ExtensionTok{scrcpy} \AttributeTok{{-}{-}crop}\OperatorTok{=}\NormalTok{1080:960:0:0    }\CommentTok{\# Para un dispositivo 1080x1920}
\end{Highlighting}
\end{Shaded}

\subsection{Bufferización}\label{bufferizaciuxf3n}

\begin{Shaded}
\begin{Highlighting}[]
\CommentTok{\# Agregar buffer para compensar fluctuaciones de red (WiFi)}
\ExtensionTok{scrcpy} \AttributeTok{{-}{-}video{-}buffer}\OperatorTok{=}\NormalTok{50    }\CommentTok{\# 50ms de buffer}
\ExtensionTok{scrcpy} \AttributeTok{{-}{-}video{-}buffer}\OperatorTok{=}\NormalTok{200   }\CommentTok{\# 200ms (más suave, más latencia)}

\CommentTok{\# Sin buffer (por defecto, menor latencia)}
\ExtensionTok{scrcpy}
\end{Highlighting}
\end{Shaded}

\subsection{Sin Reproducción de
Video}\label{sin-reproducciuxf3n-de-video}

\begin{Shaded}
\begin{Highlighting}[]
\CommentTok{\# Solo grabar, sin mostrar (útil para grabaciones largas)}
\ExtensionTok{scrcpy} \AttributeTok{{-}{-}no{-}playback} \AttributeTok{{-}{-}record}\OperatorTok{=}\NormalTok{archivo.mp4}

\CommentTok{\# O solo desactivar reproducción de video}
\ExtensionTok{scrcpy} \AttributeTok{{-}{-}no{-}video{-}playback} \AttributeTok{{-}{-}record}\OperatorTok{=}\NormalTok{archivo.mp4}
\end{Highlighting}
\end{Shaded}

\subsection{Pantalla Virtual}\label{pantalla-virtual}

\begin{Shaded}
\begin{Highlighting}[]
\CommentTok{\# Crear nueva pantalla virtual (no duplica la principal)}
\ExtensionTok{scrcpy} \AttributeTok{{-}{-}new{-}display}\OperatorTok{=}\NormalTok{1920x1080}

\CommentTok{\# Con DPI personalizado}
\ExtensionTok{scrcpy} \AttributeTok{{-}{-}new{-}display}\OperatorTok{=}\NormalTok{1920x1080/420}

\CommentTok{\# Iniciar app en pantalla virtual}
\ExtensionTok{scrcpy} \AttributeTok{{-}{-}new{-}display}\OperatorTok{=}\NormalTok{1920x1080 }\AttributeTok{{-}{-}start{-}app}\OperatorTok{=}\NormalTok{org.videolan.vlc}

\CommentTok{\# Sin decoraciones del sistema}
\ExtensionTok{scrcpy} \AttributeTok{{-}{-}new{-}display} \AttributeTok{{-}{-}no{-}vd{-}system{-}decorations} \AttributeTok{{-}{-}start{-}app}\OperatorTok{=}\NormalTok{org.mozilla.firefox}
\end{Highlighting}
\end{Shaded}

\section{Control de Audio}\label{control-de-audio}

\subsection{Habilitar/Deshabilitar
Audio}\label{habilitardeshabilitar-audio}

\begin{Shaded}
\begin{Highlighting}[]
\CommentTok{\# Sin audio (por defecto está habilitado)}
\ExtensionTok{scrcpy} \AttributeTok{{-}{-}no{-}audio}

\CommentTok{\# Solo audio, sin video}
\ExtensionTok{scrcpy} \AttributeTok{{-}{-}no{-}video} \AttributeTok{{-}{-}no{-}control}

\CommentTok{\# Audio sin ventana}
\ExtensionTok{scrcpy} \AttributeTok{{-}{-}no{-}window}  \CommentTok{\# implica {-}{-}no{-}video {-}{-}no{-}control}
\end{Highlighting}
\end{Shaded}

\subsection{Fuente de Audio}\label{fuente-de-audio}

\begin{Shaded}
\begin{Highlighting}[]
\CommentTok{\# Salida de audio del dispositivo (por defecto)}
\ExtensionTok{scrcpy} \AttributeTok{{-}{-}audio{-}source}\OperatorTok{=}\NormalTok{output}

\CommentTok{\# Micrófono del dispositivo}
\ExtensionTok{scrcpy} \AttributeTok{{-}{-}audio{-}source}\OperatorTok{=}\NormalTok{mic}

\CommentTok{\# Otras fuentes:}
\ExtensionTok{scrcpy} \AttributeTok{{-}{-}audio{-}source}\OperatorTok{=}\NormalTok{playback            }\CommentTok{\# Reproducción (apps pueden optar por no capturar)}
\ExtensionTok{scrcpy} \AttributeTok{{-}{-}audio{-}source}\OperatorTok{=}\NormalTok{mic{-}unprocessed     }\CommentTok{\# Micrófono sin procesar}
\ExtensionTok{scrcpy} \AttributeTok{{-}{-}audio{-}source}\OperatorTok{=}\NormalTok{mic{-}camcorder       }\CommentTok{\# Optimizado para video}
\ExtensionTok{scrcpy} \AttributeTok{{-}{-}audio{-}source}\OperatorTok{=}\NormalTok{mic{-}voice{-}recognition}
\ExtensionTok{scrcpy} \AttributeTok{{-}{-}audio{-}source}\OperatorTok{=}\NormalTok{mic{-}voice{-}communication}
\ExtensionTok{scrcpy} \AttributeTok{{-}{-}audio{-}source}\OperatorTok{=}\NormalTok{voice{-}call}
\ExtensionTok{scrcpy} \AttributeTok{{-}{-}audio{-}source}\OperatorTok{=}\NormalTok{voice{-}call{-}uplink}
\ExtensionTok{scrcpy} \AttributeTok{{-}{-}audio{-}source}\OperatorTok{=}\NormalTok{voice{-}call{-}downlink}
\end{Highlighting}
\end{Shaded}

\textbf{Ejemplo: Dictáfono (grabar micrófono directamente):}

\begin{Shaded}
\begin{Highlighting}[]
\ExtensionTok{scrcpy} \AttributeTok{{-}{-}audio{-}source}\OperatorTok{=}\NormalTok{mic }\AttributeTok{{-}{-}no{-}video} \AttributeTok{{-}{-}no{-}playback} \AttributeTok{{-}{-}record}\OperatorTok{=}\NormalTok{grabacion.opus}
\end{Highlighting}
\end{Shaded}

\subsection{Duplicación de Audio (Android
13+)}\label{duplicaciuxf3n-de-audio-android-13}

\begin{Shaded}
\begin{Highlighting}[]
\CommentTok{\# Reproducir audio en dispositivo Y en la computadora}
\ExtensionTok{scrcpy} \AttributeTok{{-}{-}audio{-}dup}

\CommentTok{\# Equivalente a:}
\ExtensionTok{scrcpy} \AttributeTok{{-}{-}audio{-}source}\OperatorTok{=}\NormalTok{playback }\AttributeTok{{-}{-}audio{-}dup}
\end{Highlighting}
\end{Shaded}

\subsection{Códec de Audio}\label{cuxf3dec-de-audio}

\begin{Shaded}
\begin{Highlighting}[]
\CommentTok{\# Seleccionar códec}
\ExtensionTok{scrcpy} \AttributeTok{{-}{-}audio{-}codec}\OperatorTok{=}\NormalTok{opus  }\CommentTok{\# Por defecto}
\ExtensionTok{scrcpy} \AttributeTok{{-}{-}audio{-}codec}\OperatorTok{=}\NormalTok{aac}
\ExtensionTok{scrcpy} \AttributeTok{{-}{-}audio{-}codec}\OperatorTok{=}\NormalTok{flac}
\ExtensionTok{scrcpy} \AttributeTok{{-}{-}audio{-}codec}\OperatorTok{=}\NormalTok{raw   }\CommentTok{\# PCM 16{-}bit sin comprimir}

\CommentTok{\# Ejemplo: grabar audio en FLAC de alta calidad}
\ExtensionTok{scrcpy} \AttributeTok{{-}{-}no{-}video} \AttributeTok{{-}{-}audio{-}codec}\OperatorTok{=}\NormalTok{flac }\AttributeTok{{-}{-}record}\OperatorTok{=}\NormalTok{audio.flac}

\CommentTok{\# Opciones adicionales de códec}
\ExtensionTok{scrcpy} \AttributeTok{{-}{-}audio{-}codec}\OperatorTok{=}\NormalTok{flac }\AttributeTok{{-}{-}audio{-}codec{-}options}\OperatorTok{=}\NormalTok{flac{-}compression{-}level=8}
\end{Highlighting}
\end{Shaded}

\subsection{Bitrate de Audio}\label{bitrate-de-audio}

\begin{Shaded}
\begin{Highlighting}[]
\CommentTok{\# Cambiar bitrate (por defecto 128 Kbps)}
\ExtensionTok{scrcpy} \AttributeTok{{-}{-}audio{-}bit{-}rate}\OperatorTok{=}\NormalTok{64K}
\ExtensionTok{scrcpy} \AttributeTok{{-}{-}audio{-}bit{-}rate}\OperatorTok{=}\NormalTok{64000  }\CommentTok{\# equivalente}
\ExtensionTok{scrcpy} \AttributeTok{{-}{-}audio{-}bit{-}rate}\OperatorTok{=}\NormalTok{320K   }\CommentTok{\# Alta calidad}
\end{Highlighting}
\end{Shaded}

\subsection{Bufferización de Audio}\label{bufferizaciuxf3n-de-audio}

\begin{Shaded}
\begin{Highlighting}[]
\CommentTok{\# Ajustar buffer de audio (por defecto 50ms)}
\ExtensionTok{scrcpy} \AttributeTok{{-}{-}audio{-}buffer}\OperatorTok{=}\NormalTok{40    }\CommentTok{\# Menor latencia}
\ExtensionTok{scrcpy} \AttributeTok{{-}{-}audio{-}buffer}\OperatorTok{=}\NormalTok{100   }\CommentTok{\# Más estable}
\ExtensionTok{scrcpy} \AttributeTok{{-}{-}audio{-}buffer}\OperatorTok{=}\NormalTok{200   }\CommentTok{\# Muy estable, más latencia}

\CommentTok{\# Buffer de salida de audio (por defecto 5ms)}
\CommentTok{\# Solo cambiar si hay audio robótico/entrecortado}
\ExtensionTok{scrcpy} \AttributeTok{{-}{-}audio{-}output{-}buffer}\OperatorTok{=}\NormalTok{10}
\end{Highlighting}
\end{Shaded}

\section{Grabación de Pantalla}\label{grabaciuxf3n-de-pantalla}

\subsection{Grabación Básica}\label{grabaciuxf3n-buxe1sica}

\begin{Shaded}
\begin{Highlighting}[]
\CommentTok{\# Grabar video y audio}
\ExtensionTok{scrcpy} \AttributeTok{{-}{-}record}\OperatorTok{=}\NormalTok{archivo.mp4}
\ExtensionTok{scrcpy} \AttributeTok{{-}r}\NormalTok{ archivo.mkv  }\CommentTok{\# versión corta}

\CommentTok{\# Solo video}
\ExtensionTok{scrcpy} \AttributeTok{{-}{-}no{-}audio} \AttributeTok{{-}{-}record}\OperatorTok{=}\NormalTok{video.mp4}

\CommentTok{\# Solo audio}
\ExtensionTok{scrcpy} \AttributeTok{{-}{-}no{-}video} \AttributeTok{{-}{-}record}\OperatorTok{=}\NormalTok{audio.opus}
\ExtensionTok{scrcpy} \AttributeTok{{-}{-}no{-}video} \AttributeTok{{-}{-}audio{-}codec}\OperatorTok{=}\NormalTok{aac }\AttributeTok{{-}{-}record}\OperatorTok{=}\NormalTok{audio.aac}
\ExtensionTok{scrcpy} \AttributeTok{{-}{-}no{-}video} \AttributeTok{{-}{-}audio{-}codec}\OperatorTok{=}\NormalTok{flac }\AttributeTok{{-}{-}record}\OperatorTok{=}\NormalTok{audio.flac}
\ExtensionTok{scrcpy} \AttributeTok{{-}{-}no{-}video} \AttributeTok{{-}{-}audio{-}codec}\OperatorTok{=}\NormalTok{raw }\AttributeTok{{-}{-}record}\OperatorTok{=}\NormalTok{audio.wav}
\end{Highlighting}
\end{Shaded}

\subsection{Formatos Disponibles}\label{formatos-disponibles}

{\def\LTcaptype{none} % do not increment counter
\begin{longtable}[]{@{}lll@{}}
\toprule\noalign{}
Formato & Extensiones & Uso \\
\midrule\noalign{}
\endhead
\bottomrule\noalign{}
\endlastfoot
MP4 & \texttt{.mp4}, \texttt{.m4a}, \texttt{.aac} & Video y audio,
compatible universalmente \\
Matroska & \texttt{.mkv}, \texttt{.mka} & Video y audio, mejor para
edición \\
OPUS & \texttt{.opus} & Solo audio, excelente compresión \\
FLAC & \texttt{.flac} & Solo audio, sin pérdida \\
WAV & \texttt{.wav} & Solo audio, sin comprimir \\
\end{longtable}
}

\textbf{Especificar formato explícitamente:}

\begin{Shaded}
\begin{Highlighting}[]
\ExtensionTok{scrcpy} \AttributeTok{{-}{-}record}\OperatorTok{=}\NormalTok{archivo }\AttributeTok{{-}{-}record{-}format}\OperatorTok{=}\NormalTok{mkv}
\end{Highlighting}
\end{Shaded}

\subsection{Grabación Sin
Reproducción}\label{grabaciuxf3n-sin-reproducciuxf3n}

\begin{Shaded}
\begin{Highlighting}[]
\CommentTok{\# Grabar sin mostrar en pantalla (ahorra recursos)}
\ExtensionTok{scrcpy} \AttributeTok{{-}{-}no{-}playback} \AttributeTok{{-}{-}record}\OperatorTok{=}\NormalTok{archivo.mp4}

\CommentTok{\# Sin ventana (modo background)}
\ExtensionTok{scrcpy} \AttributeTok{{-}{-}no{-}window} \AttributeTok{{-}{-}record}\OperatorTok{=}\NormalTok{archivo.mp4}
\CommentTok{\# Interrumpir con Ctrl+C}

\CommentTok{\# Solo grabar video, seguir escuchando audio}
\ExtensionTok{scrcpy} \AttributeTok{{-}{-}no{-}video{-}playback} \AttributeTok{{-}{-}record}\OperatorTok{=}\NormalTok{archivo.mp4}
\end{Highlighting}
\end{Shaded}

\subsection{Límite de Tiempo}\label{luxedmite-de-tiempo}

\begin{Shaded}
\begin{Highlighting}[]
\CommentTok{\# Grabar por 20 segundos}
\ExtensionTok{scrcpy} \AttributeTok{{-}{-}record}\OperatorTok{=}\NormalTok{archivo.mp4 }\AttributeTok{{-}{-}time{-}limit}\OperatorTok{=}\NormalTok{20}

\CommentTok{\# También funciona sin grabar}
\ExtensionTok{scrcpy} \AttributeTok{{-}{-}time{-}limit}\OperatorTok{=}\NormalTok{30  }\CommentTok{\# Duplicar por 30 segundos y cerrar}
\end{Highlighting}
\end{Shaded}

\subsection{Ejemplos de Grabación
Profesional}\label{ejemplos-de-grabaciuxf3n-profesional}

\subsubsection{Tutorial de App Móvil (Alta
Calidad)}\label{tutorial-de-app-muxf3vil-alta-calidad}

\begin{Shaded}
\begin{Highlighting}[]
\ExtensionTok{scrcpy} \AttributeTok{{-}{-}video{-}codec}\OperatorTok{=}\NormalTok{h265 }\AttributeTok{{-}m1920} \AttributeTok{{-}{-}max{-}fps}\OperatorTok{=}\NormalTok{60 }\DataTypeTok{\textbackslash{}}
  \AttributeTok{{-}{-}audio{-}codec}\OperatorTok{=}\NormalTok{aac }\DataTypeTok{\textbackslash{}}
  \AttributeTok{{-}{-}record}\OperatorTok{=}\NormalTok{tutorial{-}app.mp4}
\end{Highlighting}
\end{Shaded}

\subsubsection{Presentación Académica
(Balance)}\label{presentaciuxf3n-acaduxe9mica-balance}

\begin{Shaded}
\begin{Highlighting}[]
\ExtensionTok{scrcpy} \AttributeTok{{-}m1920} \AttributeTok{{-}{-}max{-}fps}\OperatorTok{=}\NormalTok{30 }\DataTypeTok{\textbackslash{}}
  \AttributeTok{{-}{-}video{-}bit{-}rate}\OperatorTok{=}\NormalTok{4M }\DataTypeTok{\textbackslash{}}
  \AttributeTok{{-}{-}audio{-}bit{-}rate}\OperatorTok{=}\NormalTok{128K }\DataTypeTok{\textbackslash{}}
  \AttributeTok{{-}{-}record}\OperatorTok{=}\NormalTok{presentacion.mkv}
\end{Highlighting}
\end{Shaded}

\subsubsection{Grabación de Podcast (Solo
Audio)}\label{grabaciuxf3n-de-podcast-solo-audio}

\begin{Shaded}
\begin{Highlighting}[]
\ExtensionTok{scrcpy} \AttributeTok{{-}{-}audio{-}source}\OperatorTok{=}\NormalTok{mic }\DataTypeTok{\textbackslash{}}
  \AttributeTok{{-}{-}audio{-}codec}\OperatorTok{=}\NormalTok{opus }\DataTypeTok{\textbackslash{}}
  \AttributeTok{{-}{-}no{-}video} \DataTypeTok{\textbackslash{}}
  \AttributeTok{{-}{-}no{-}playback} \DataTypeTok{\textbackslash{}}
  \AttributeTok{{-}{-}record}\OperatorTok{=}\NormalTok{podcast.opus}
\end{Highlighting}
\end{Shaded}

\subsubsection{Demo de Análisis Econométrico (Pantalla +
Audio)}\label{demo-de-anuxe1lisis-economuxe9trico-pantalla-audio}

\begin{Shaded}
\begin{Highlighting}[]
\ExtensionTok{scrcpy} \AttributeTok{{-}{-}video{-}codec}\OperatorTok{=}\NormalTok{h264 }\DataTypeTok{\textbackslash{}}
  \AttributeTok{{-}m1920} \DataTypeTok{\textbackslash{}}
  \AttributeTok{{-}{-}max{-}fps}\OperatorTok{=}\NormalTok{60 }\DataTypeTok{\textbackslash{}}
  \AttributeTok{{-}{-}stay{-}awake} \DataTypeTok{\textbackslash{}}
  \AttributeTok{{-}{-}record}\OperatorTok{=}\NormalTok{demo{-}econometria.mp4}
\end{Highlighting}
\end{Shaded}

\begin{center}\rule{0.5\linewidth}{0.5pt}\end{center}

\section{Atajos de Teclado}\label{atajos-de-teclado}

\textbf{MOD} es el modificador de atajos. Por defecto: \texttt{Alt}
(izquierda) o \texttt{Super} (Windows/Cmd).

\textbf{Cambiar MOD:}

\begin{Shaded}
\begin{Highlighting}[]
\CommentTok{\# Usar Ctrl derecho}
\ExtensionTok{scrcpy} \AttributeTok{{-}{-}shortcut{-}mod}\OperatorTok{=}\NormalTok{rctrl}

\CommentTok{\# Usar múltiples modificadores}
\ExtensionTok{scrcpy} \AttributeTok{{-}{-}shortcut{-}mod}\OperatorTok{=}\NormalTok{lctrl,lsuper}
\end{Highlighting}
\end{Shaded}

\subsection{Tabla de Atajos}\label{tabla-de-atajos}

{\def\LTcaptype{none} % do not increment counter
\begin{longtable}[]{@{}ll@{}}
\toprule\noalign{}
Acción & Atajo \\
\midrule\noalign{}
\endhead
\bottomrule\noalign{}
\endlastfoot
\textbf{Ventana} & \\
Pantalla completa & \texttt{MOD+f} \\
Redimensionar 1:1 (pixel-perfect) & \texttt{MOD+g} \\
Quitar bordes negros & \texttt{MOD+w} o doble-clic \\
Siempre visible & Usar \texttt{-\/-always-on-top} \\
\textbf{Rotación/Transformación} & \\
Rotar a la izquierda & \texttt{MOD+←} \\
Rotar a la derecha & \texttt{MOD+→} \\
Voltear horizontalmente & \texttt{MOD+Shift+←} o \texttt{MOD+Shift+→} \\
Voltear verticalmente & \texttt{MOD+Shift+↑} o \texttt{MOD+Shift+↓} \\
Rotar dispositivo & \texttt{MOD+r} \\
\textbf{Pantalla del Dispositivo} & \\
Apagar pantalla & \texttt{MOD+o} \\
Encender pantalla & \texttt{MOD+Shift+o} \\
\textbf{Botones Android} & \\
HOME & \texttt{MOD+h} o clic-medio \\
BACK & \texttt{MOD+b} o \texttt{MOD+Backspace} o clic-derecho \\
APP\_SWITCH (recientes) & \texttt{MOD+s} o 4to botón mouse \\
MENU & \texttt{MOD+m} \\
\textbf{Volumen} & \\
Subir volumen & \texttt{MOD+↑} \\
Bajar volumen & \texttt{MOD+↓} \\
Encender/apagar & \texttt{MOD+p} \\
\textbf{Paneles} & \\
Notificaciones & \texttt{MOD+n} o 5to botón mouse \\
Configuración & \texttt{MOD+n+n} (doble) \\
Colapsar paneles & \texttt{MOD+Shift+n} \\
\textbf{Copiar/Pegar} & \\
Copiar & \texttt{MOD+c} \\
Cortar & \texttt{MOD+x} \\
Pegar & \texttt{MOD+v} \\
Pegar como eventos & \texttt{MOD+Shift+v} \\
\textbf{Control} & \\
Configuración teclado (HID) & \texttt{MOD+k} \\
Mostrar/ocultar FPS & \texttt{MOD+i} \\
Pausar duplicación & \texttt{MOD+z} \\
Reanudar & \texttt{MOD+Shift+z} \\
\textbf{Gestos} & \\
Pinch-to-zoom & \texttt{Ctrl+clic+mover} \\
Inclinación vertical & \texttt{Shift+clic+mover} \\
Inclinación horizontal & \texttt{Ctrl+Shift+clic+mover} \\
\textbf{Archivos} & \\
Instalar APK & Arrastrar archivo \texttt{.apk} \\
Enviar archivo & Arrastrar archivo (no-APK) \\
\end{longtable}
}

\begin{center}\rule{0.5\linewidth}{0.5pt}\end{center}

\section{Casos de Uso Profesionales}\label{casos-de-uso-profesionales}

\subsection{Caso 1: Presentación Académica con Diapositivas
Móviles}\label{caso-1-presentaciuxf3n-acaduxe9mica-con-diapositivas-muxf3viles}

\textbf{Escenario:} Presentar análisis econométrico desde app móvil en
clase.

\begin{Shaded}
\begin{Highlighting}[]
\CommentTok{\# Configuración óptima}
\ExtensionTok{scrcpy} \AttributeTok{{-}m1920} \DataTypeTok{\textbackslash{}}
  \AttributeTok{{-}{-}max{-}fps}\OperatorTok{=}\NormalTok{60 }\DataTypeTok{\textbackslash{}}
  \AttributeTok{{-}{-}video{-}bit{-}rate}\OperatorTok{=}\NormalTok{8M }\DataTypeTok{\textbackslash{}}
  \AttributeTok{{-}{-}no{-}audio} \DataTypeTok{\textbackslash{}}
  \AttributeTok{{-}{-}always{-}on{-}top} \DataTypeTok{\textbackslash{}}
  \AttributeTok{{-}{-}stay{-}awake} \DataTypeTok{\textbackslash{}}
  \AttributeTok{{-}f}
\end{Highlighting}
\end{Shaded}

\textbf{Explicación:}

\begin{itemize}
\tightlist
\item
  \texttt{-m1920}: Resolución óptima para proyector
\item
  \texttt{-\/-max-fps=60}: Fluidez en animaciones
\item
  \texttt{-\/-no-audio}: Sin distracciones
\item
  \texttt{-\/-always-on-top}: Siempre visible
\item
  \texttt{-\/-stay-awake}: Dispositivo no se apaga
\item
  \texttt{-f}: Pantalla completa
\end{itemize}

\subsection{Caso 2: Grabar Tutorial de App
Econométrica}\label{caso-2-grabar-tutorial-de-app-economuxe9trica}

\textbf{Escenario:} Grabar tutorial de uso de app de estadística.

\begin{Shaded}
\begin{Highlighting}[]
\CommentTok{\# Grabación profesional}
\ExtensionTok{scrcpy} \AttributeTok{{-}{-}video{-}codec}\OperatorTok{=}\NormalTok{h265 }\DataTypeTok{\textbackslash{}}
  \AttributeTok{{-}m1920} \DataTypeTok{\textbackslash{}}
  \AttributeTok{{-}{-}max{-}fps}\OperatorTok{=}\NormalTok{60 }\DataTypeTok{\textbackslash{}}
  \AttributeTok{{-}{-}video{-}bit{-}rate}\OperatorTok{=}\NormalTok{8M }\DataTypeTok{\textbackslash{}}
  \AttributeTok{{-}{-}audio{-}source}\OperatorTok{=}\NormalTok{mic }\DataTypeTok{\textbackslash{}}
  \AttributeTok{{-}{-}audio{-}codec}\OperatorTok{=}\NormalTok{aac }\DataTypeTok{\textbackslash{}}
  \AttributeTok{{-}{-}audio{-}bit{-}rate}\OperatorTok{=}\NormalTok{192K }\DataTypeTok{\textbackslash{}}
  \AttributeTok{{-}{-}record}\OperatorTok{=}\NormalTok{tutorial{-}app{-}estadistica.mp4 }\DataTypeTok{\textbackslash{}}
  \AttributeTok{{-}{-}stay{-}awake} \DataTypeTok{\textbackslash{}}
  \AttributeTok{{-}{-}show{-}touches}
\end{Highlighting}
\end{Shaded}

\textbf{Explicación:}

\begin{itemize}
\tightlist
\item
  \texttt{-\/-video-codec=h265}: Mejor calidad/compresión
\item
  \texttt{-\/-audio-source=mic}: Narración con tu voz
\item
  \texttt{-\/-show-touches}: Muestra dónde tocas
\item
  \texttt{-\/-stay-awake}: No se apaga durante grabación
\end{itemize}

\subsection{Caso 3: Demostración en Vivo con
Webcam}\label{caso-3-demostraciuxf3n-en-vivo-con-webcam}

\textbf{Escenario:} Usar cámara del dispositivo como webcam para
videollamada.

\begin{Shaded}
\begin{Highlighting}[]
\CommentTok{\# Configurar como webcam (Linux)}
\ExtensionTok{scrcpy} \AttributeTok{{-}{-}video{-}source}\OperatorTok{=}\NormalTok{camera }\DataTypeTok{\textbackslash{}}
  \AttributeTok{{-}{-}camera{-}size}\OperatorTok{=}\NormalTok{1920x1080 }\DataTypeTok{\textbackslash{}}
  \AttributeTok{{-}{-}camera{-}facing}\OperatorTok{=}\NormalTok{front }\DataTypeTok{\textbackslash{}}
  \AttributeTok{{-}{-}v4l2{-}sink}\OperatorTok{=}\NormalTok{/dev/video2 }\DataTypeTok{\textbackslash{}}
  \AttributeTok{{-}{-}no{-}playback}
\end{Highlighting}
\end{Shaded}

\textbf{Luego usar \texttt{/dev/video2} en Zoom, Google Meet, etc.}

\subsection{Caso 4: Control Remoto para
Demostraciones}\label{caso-4-control-remoto-para-demostraciones}

\textbf{Escenario:} Controlar dispositivo de forma remota vía WiFi.

\begin{Shaded}
\begin{Highlighting}[]
\CommentTok{\# 1. Conectar por USB una vez}
\ExtensionTok{scrcpy} \AttributeTok{{-}{-}tcpip}

\CommentTok{\# 2. Desconectar cable}
\CommentTok{\# 3. Ahora funciona inalámbrico}
\ExtensionTok{scrcpy} \AttributeTok{{-}m1024} \AttributeTok{{-}{-}max{-}fps}\OperatorTok{=}\NormalTok{30}
\end{Highlighting}
\end{Shaded}

\subsection{Caso 5: Documentación de Aplicación con
Capturas}\label{caso-5-documentaciuxf3n-de-aplicaciuxf3n-con-capturas}

\textbf{Escenario:} Documentar flujo de trabajo de app.

\begin{Shaded}
\begin{Highlighting}[]
\CommentTok{\# Grabar sesión completa}
\ExtensionTok{scrcpy} \AttributeTok{{-}{-}video{-}codec}\OperatorTok{=}\NormalTok{h265 }\DataTypeTok{\textbackslash{}}
  \AttributeTok{{-}m1920} \DataTypeTok{\textbackslash{}}
  \AttributeTok{{-}{-}max{-}fps}\OperatorTok{=}\NormalTok{30 }\DataTypeTok{\textbackslash{}}
  \AttributeTok{{-}{-}record}\OperatorTok{=}\NormalTok{documentacion{-}flujo.mkv }\DataTypeTok{\textbackslash{}}
  \AttributeTok{{-}{-}no{-}audio} \DataTypeTok{\textbackslash{}}
  \AttributeTok{{-}{-}show{-}touches}

\CommentTok{\# Luego extraer frames con ffmpeg:}
\ExtensionTok{ffmpeg} \AttributeTok{{-}i}\NormalTok{ documentacion{-}flujo.mkv }\AttributeTok{{-}vf}\NormalTok{ fps=1/5 captura\_\%04d.png}
\end{Highlighting}
\end{Shaded}

\subsection{Caso 6: Podcast/Entrevista con Audio
Móvil}\label{caso-6-podcastentrevista-con-audio-muxf3vil}

\textbf{Escenario:} Grabar entrevista usando el teléfono como grabadora.

\begin{Shaded}
\begin{Highlighting}[]
\ExtensionTok{scrcpy} \AttributeTok{{-}{-}audio{-}source}\OperatorTok{=}\NormalTok{mic }\DataTypeTok{\textbackslash{}}
  \AttributeTok{{-}{-}audio{-}codec}\OperatorTok{=}\NormalTok{flac }\DataTypeTok{\textbackslash{}}
  \AttributeTok{{-}{-}audio{-}buffer}\OperatorTok{=}\NormalTok{200 }\DataTypeTok{\textbackslash{}}
  \AttributeTok{{-}{-}no{-}video} \DataTypeTok{\textbackslash{}}
  \AttributeTok{{-}{-}no{-}window} \DataTypeTok{\textbackslash{}}
  \AttributeTok{{-}{-}record}\OperatorTok{=}\NormalTok{entrevista{-}}\VariableTok{$(}\FunctionTok{date}\NormalTok{ +\%Y\%m\%d}\VariableTok{)}\NormalTok{.flac}
\end{Highlighting}
\end{Shaded}

\begin{center}\rule{0.5\linewidth}{0.5pt}\end{center}

\section{Troubleshooting}\label{troubleshooting}

\subsection{Problema 1: ``No se detecta el
dispositivo''}\label{problema-1-no-se-detecta-el-dispositivo}

\textbf{Síntomas:}

\begin{verbatim}
adb devices
# List of devices attached
# (vacío)
\end{verbatim}

\textbf{Soluciones:}

\begin{Shaded}
\begin{Highlighting}[]
\CommentTok{\# 1. Verificar cable USB (usar cable de datos, no solo carga)}

\CommentTok{\# 2. Verificar depuración USB en dispositivo}
\CommentTok{\# Ajustes → Opciones de desarrollador → Depuración USB}

\CommentTok{\# 3. Autorizar computadora en dispositivo (popup)}

\CommentTok{\# 4. Reiniciar servidor adb}
\ExtensionTok{adb}\NormalTok{ kill{-}server}
\ExtensionTok{adb}\NormalTok{ start{-}server}
\ExtensionTok{adb}\NormalTok{ devices}

\CommentTok{\# 5. Verificar permisos}
\FunctionTok{groups} \KeywordTok{|} \FunctionTok{grep}\NormalTok{ adbusers}
\CommentTok{\# Si no estás: sudo usermod {-}aG adbusers $USER}

\CommentTok{\# 6. Verificar reglas udev}
\FunctionTok{ls}\NormalTok{ /etc/udev/rules.d/51{-}android.rules}
\CommentTok{\# Si no existe, crearla (ver sección Configuración Inicial)}

\CommentTok{\# 7. Probar con sudo (temporal, para descartar permisos)}
\FunctionTok{sudo}\NormalTok{ adb devices}
\end{Highlighting}
\end{Shaded}

\subsection{Problema 2: ``Error: could not install
server''}\label{problema-2-error-could-not-install-server}

\textbf{Causa:} No hay espacio en \texttt{/data/local/tmp} o permisos
incorrectos.

\textbf{Soluciones:}

\begin{Shaded}
\begin{Highlighting}[]
\CommentTok{\# 1. Limpiar archivos temporales en dispositivo}
\ExtensionTok{adb}\NormalTok{ shell rm }\AttributeTok{{-}rf}\NormalTok{ /data/local/tmp/scrcpy{-}server}\PreprocessorTok{*}

\CommentTok{\# 2. Verificar espacio}
\ExtensionTok{adb}\NormalTok{ shell df }\AttributeTok{{-}h}\NormalTok{ /data}

\CommentTok{\# 3. Si falta espacio, limpiar caché}
\CommentTok{\# En dispositivo: Ajustes → Almacenamiento → Limpiar caché}

\CommentTok{\# 4. Intentar de nuevo}
\ExtensionTok{scrcpy}
\end{Highlighting}
\end{Shaded}

\subsection{Problema 3: ``Pantalla negra pero audio
funciona''}\label{problema-3-pantalla-negra-pero-audio-funciona}

\textbf{Causas posibles:}

\begin{itemize}
\tightlist
\item
  Códec de video no soportado
\item
  Resolución muy alta
\item
  Problema de decodificación
\end{itemize}

\textbf{Soluciones:}

\begin{Shaded}
\begin{Highlighting}[]
\CommentTok{\# 1. Reducir resolución}
\ExtensionTok{scrcpy} \AttributeTok{{-}m1024}

\CommentTok{\# 2. Cambiar códec}
\ExtensionTok{scrcpy} \AttributeTok{{-}{-}video{-}codec}\OperatorTok{=}\NormalTok{h264  }\CommentTok{\# Más compatible}

\CommentTok{\# 3. Listar codificadores disponibles}
\ExtensionTok{scrcpy} \AttributeTok{{-}{-}list{-}encoders}

\CommentTok{\# 4. Probar con otro codificador}
\ExtensionTok{scrcpy} \AttributeTok{{-}{-}video{-}codec}\OperatorTok{=}\NormalTok{h264 }\AttributeTok{{-}{-}video{-}encoder}\OperatorTok{=}\NormalTok{OMX.qcom.video.encoder.avc}

\CommentTok{\# 5. Verificar FFmpeg}
\ExtensionTok{ffmpeg} \AttributeTok{{-}version}
\end{Highlighting}
\end{Shaded}

\subsection{Problema 4: ``Audio entrecortado o
robótico''}\label{problema-4-audio-entrecortado-o-robuxf3tico}

\textbf{Soluciones:}

\begin{Shaded}
\begin{Highlighting}[]
\CommentTok{\# 1. Aumentar buffer de audio}
\ExtensionTok{scrcpy} \AttributeTok{{-}{-}audio{-}buffer}\OperatorTok{=}\NormalTok{200}

\CommentTok{\# 2. Aumentar buffer de salida}
\ExtensionTok{scrcpy} \AttributeTok{{-}{-}audio{-}output{-}buffer}\OperatorTok{=}\NormalTok{10}

\CommentTok{\# 3. Cambiar códec}
\ExtensionTok{scrcpy} \AttributeTok{{-}{-}audio{-}codec}\OperatorTok{=}\NormalTok{aac  }\CommentTok{\# En lugar de opus}

\CommentTok{\# 4. Reducir calidad de video (libera CPU)}
\ExtensionTok{scrcpy} \AttributeTok{{-}m1024} \AttributeTok{{-}{-}max{-}fps}\OperatorTok{=}\NormalTok{30}

\CommentTok{\# 5. Conexión por cable (en lugar de WiFi)}
\end{Highlighting}
\end{Shaded}

\subsection{Problema 5: ``Latencia alta en
WiFi''}\label{problema-5-latencia-alta-en-wifi}

\textbf{Soluciones:}

\begin{Shaded}
\begin{Highlighting}[]
\CommentTok{\# 1. Usar cable USB}
\CommentTok{\# WiFi siempre tendrá más latencia que USB}

\CommentTok{\# 2. Reducir calidad}
\ExtensionTok{scrcpy} \AttributeTok{{-}m1024} \AttributeTok{{-}{-}max{-}fps}\OperatorTok{=}\NormalTok{30 }\AttributeTok{{-}{-}video{-}bit{-}rate}\OperatorTok{=}\NormalTok{2M}

\CommentTok{\# 3. Agregar buffer}
\ExtensionTok{scrcpy} \AttributeTok{{-}{-}video{-}buffer}\OperatorTok{=}\NormalTok{100 }\AttributeTok{{-}{-}audio{-}buffer}\OperatorTok{=}\NormalTok{200}

\CommentTok{\# 4. Verificar red WiFi (2.4 GHz vs 5 GHz)}
\CommentTok{\# 5 GHz tiene menor latencia pero menor alcance}

\CommentTok{\# 5. Reducir interferencias (cerrar otras apps de red)}
\end{Highlighting}
\end{Shaded}

\subsection{Problema 6: ``No se puede
copiar/pegar''}\label{problema-6-no-se-puede-copiarpegar}

\textbf{Requisito:} Android 7 o superior.

\textbf{Soluciones:}

\begin{Shaded}
\begin{Highlighting}[]
\CommentTok{\# 1. Verificar versión de Android}
\ExtensionTok{adb}\NormalTok{ shell getprop ro.build.version.release}

\CommentTok{\# 2. Usar atajos correctos}
\CommentTok{\# MOD+c (copiar), MOD+v (pegar)}

\CommentTok{\# 3. Si no funciona, usar modo legacy}
\ExtensionTok{scrcpy} \AttributeTok{{-}{-}legacy{-}paste}

\CommentTok{\# 4. Pegar como eventos de teclas}
\CommentTok{\# MOD+Shift+v}
\end{Highlighting}
\end{Shaded}

\subsection{Problema 7: ``Dispositivo Xiaomi no
responde''}\label{problema-7-dispositivo-xiaomi-no-responde}

\textbf{Solución específica Xiaomi:}

\begin{Shaded}
\begin{Highlighting}[]
\CommentTok{\# Habilitar opción adicional en Opciones de desarrollador:}
\CommentTok{\# "Depuración USB (Configuración de seguridad)"}

\CommentTok{\# Reiniciar dispositivo}

\CommentTok{\# Verificar}
\ExtensionTok{adb}\NormalTok{ devices}
\ExtensionTok{scrcpy}
\end{Highlighting}
\end{Shaded}

\subsection{Problema 8: ``OTG no funciona en
Windows''}\label{problema-8-otg-no-funciona-en-windows}

\textbf{Causa:} Windows no permite abrir dispositivo USB si ya está
abierto por otro proceso (adb).

\textbf{Solución:}

\begin{Shaded}
\begin{Highlighting}[]
\CommentTok{\# Usar OTG requiere que adb NO esté corriendo}
\ExtensionTok{adb}\NormalTok{ kill{-}server}

\CommentTok{\# Luego usar OTG}
\ExtensionTok{scrcpy} \AttributeTok{{-}{-}otg}

\CommentTok{\# Nota: En OTG no hay duplicación de pantalla}
\end{Highlighting}
\end{Shaded}

\section{Comandos de Referencia
Rápida}\label{comandos-de-referencia-ruxe1pida}

\subsection{Uso General}\label{uso-general}

\begin{Shaded}
\begin{Highlighting}[]
\CommentTok{\# Básico}
\ExtensionTok{scrcpy}                          \CommentTok{\# Duplicar pantalla}
\ExtensionTok{scrcpy} \AttributeTok{{-}m1024}                   \CommentTok{\# Limitar a 1024px}
\ExtensionTok{scrcpy} \AttributeTok{{-}f}                       \CommentTok{\# Pantalla completa}
\ExtensionTok{scrcpy} \AttributeTok{{-}w}                       \CommentTok{\# Mantener despierto}
\ExtensionTok{scrcpy} \AttributeTok{{-}S}                       \CommentTok{\# Apagar pantalla}

\CommentTok{\# Selección de dispositivo}
\ExtensionTok{scrcpy} \AttributeTok{{-}s}\NormalTok{ 0123456789ABCDEF      }\CommentTok{\# Por serial}
\ExtensionTok{scrcpy} \AttributeTok{{-}d}                       \CommentTok{\# USB (único)}
\ExtensionTok{scrcpy} \AttributeTok{{-}e}                       \CommentTok{\# TCP/IP (único)}
\ExtensionTok{scrcpy} \AttributeTok{{-}{-}tcpip}\OperatorTok{=}\NormalTok{192.168.1.100    }\CommentTok{\# TCP/IP directo}

\CommentTok{\# Calidad}
\ExtensionTok{scrcpy} \AttributeTok{{-}m1920} \AttributeTok{{-}{-}max{-}fps}\OperatorTok{=}\NormalTok{60 }\AttributeTok{{-}b}\NormalTok{ 8M}
\ExtensionTok{scrcpy} \AttributeTok{{-}{-}video{-}codec}\OperatorTok{=}\NormalTok{h265}
\ExtensionTok{scrcpy} \AttributeTok{{-}{-}audio{-}codec}\OperatorTok{=}\NormalTok{aac}
\end{Highlighting}
\end{Shaded}

\subsection{Grabación}\label{grabaciuxf3n}

\begin{Shaded}
\begin{Highlighting}[]
\CommentTok{\# Video y audio}
\ExtensionTok{scrcpy} \AttributeTok{{-}r}\NormalTok{ archivo.mp4}
\ExtensionTok{scrcpy} \AttributeTok{{-}{-}record}\OperatorTok{=}\NormalTok{archivo.mkv}

\CommentTok{\# Solo video}
\ExtensionTok{scrcpy} \AttributeTok{{-}{-}no{-}audio} \AttributeTok{{-}r}\NormalTok{ video.mp4}

\CommentTok{\# Solo audio}
\ExtensionTok{scrcpy} \AttributeTok{{-}{-}no{-}video} \AttributeTok{{-}r}\NormalTok{ audio.opus}

\CommentTok{\# Sin mostrar}
\ExtensionTok{scrcpy} \AttributeTok{{-}{-}no{-}playback} \AttributeTok{{-}r}\NormalTok{ grabacion.mp4}

\CommentTok{\# Límite de tiempo}
\ExtensionTok{scrcpy} \AttributeTok{{-}r}\NormalTok{ archivo.mp4 }\AttributeTok{{-}{-}time{-}limit}\OperatorTok{=}\NormalTok{60}
\end{Highlighting}
\end{Shaded}

\subsection{Audio}\label{audio}

\begin{Shaded}
\begin{Highlighting}[]
\CommentTok{\# Sin audio}
\ExtensionTok{scrcpy} \AttributeTok{{-}{-}no{-}audio}

\CommentTok{\# Fuentes de audio}
\ExtensionTok{scrcpy} \AttributeTok{{-}{-}audio{-}source}\OperatorTok{=}\NormalTok{mic       }\CommentTok{\# Micrófono}
\ExtensionTok{scrcpy} \AttributeTok{{-}{-}audio{-}source}\OperatorTok{=}\NormalTok{output    }\CommentTok{\# Salida (default)}

\CommentTok{\# Solo audio}
\ExtensionTok{scrcpy} \AttributeTok{{-}{-}no{-}video} \AttributeTok{{-}{-}no{-}control}

\CommentTok{\# Duplicar audio (dispositivo + PC)}
\ExtensionTok{scrcpy} \AttributeTok{{-}{-}audio{-}dup}
\end{Highlighting}
\end{Shaded}

\subsection{Control}\label{control}

\begin{Shaded}
\begin{Highlighting}[]
\CommentTok{\# Sin control}
\ExtensionTok{scrcpy} \AttributeTok{{-}n}

\CommentTok{\# Sin video ni audio}
\ExtensionTok{scrcpy} \AttributeTok{{-}{-}no{-}video} \AttributeTok{{-}{-}no{-}audio}

\CommentTok{\# Mostrar toques}
\ExtensionTok{scrcpy} \AttributeTok{{-}{-}show{-}touches}
\ExtensionTok{scrcpy} \AttributeTok{{-}t}
\end{Highlighting}
\end{Shaded}

\subsection{Teclado/Mouse}\label{tecladomouse}

\begin{Shaded}
\begin{Highlighting}[]
\CommentTok{\# Simulación física (HID)}
\ExtensionTok{scrcpy} \AttributeTok{{-}K}                       \CommentTok{\# Teclado UHID}
\ExtensionTok{scrcpy} \AttributeTok{{-}M}                       \CommentTok{\# Mouse UHID}
\ExtensionTok{scrcpy} \AttributeTok{{-}KM}                      \CommentTok{\# Ambos}

\CommentTok{\# Modo OTG (sin adb)}
\ExtensionTok{scrcpy} \AttributeTok{{-}{-}otg}
\ExtensionTok{scrcpy} \AttributeTok{{-}{-}otg} \AttributeTok{{-}G}                 \CommentTok{\# Con gamepad}
\end{Highlighting}
\end{Shaded}

\subsection{Pantalla Virtual}\label{pantalla-virtual-1}

\begin{Shaded}
\begin{Highlighting}[]
\CommentTok{\# Crear display virtual}
\ExtensionTok{scrcpy} \AttributeTok{{-}{-}new{-}display}\OperatorTok{=}\NormalTok{1920x1080}

\CommentTok{\# Con app}
\ExtensionTok{scrcpy} \AttributeTok{{-}{-}new{-}display} \AttributeTok{{-}{-}start{-}app}\OperatorTok{=}\NormalTok{org.videolan.vlc}

\CommentTok{\# Sin decoraciones}
\ExtensionTok{scrcpy} \AttributeTok{{-}{-}new{-}display} \AttributeTok{{-}{-}no{-}vd{-}system{-}decorations}
\end{Highlighting}
\end{Shaded}

\subsection{Cámara}\label{cuxe1mara}

\begin{Shaded}
\begin{Highlighting}[]
\CommentTok{\# Duplicar cámara}
\ExtensionTok{scrcpy} \AttributeTok{{-}{-}video{-}source}\OperatorTok{=}\NormalTok{camera}

\CommentTok{\# Seleccionar cámara}
\ExtensionTok{scrcpy} \AttributeTok{{-}{-}video{-}source}\OperatorTok{=}\NormalTok{camera }\AttributeTok{{-}{-}camera{-}facing}\OperatorTok{=}\NormalTok{front}
\ExtensionTok{scrcpy} \AttributeTok{{-}{-}video{-}source}\OperatorTok{=}\NormalTok{camera }\AttributeTok{{-}{-}camera{-}id}\OperatorTok{=}\NormalTok{0}

\CommentTok{\# Webcam (Linux)}
\ExtensionTok{scrcpy} \AttributeTok{{-}{-}video{-}source}\OperatorTok{=}\NormalTok{camera }\DataTypeTok{\textbackslash{}}
  \AttributeTok{{-}{-}camera{-}size}\OperatorTok{=}\NormalTok{1920x1080 }\DataTypeTok{\textbackslash{}}
  \AttributeTok{{-}{-}v4l2{-}sink}\OperatorTok{=}\NormalTok{/dev/video2 }\DataTypeTok{\textbackslash{}}
  \AttributeTok{{-}{-}no{-}playback}
\end{Highlighting}
\end{Shaded}

\subsection{Combinaciones Útiles}\label{combinaciones-uxfatiles}

\begin{Shaded}
\begin{Highlighting}[]
\CommentTok{\# Presentación fluida}
\ExtensionTok{scrcpy} \AttributeTok{{-}m1024} \AttributeTok{{-}{-}max{-}fps}\OperatorTok{=}\NormalTok{60 }\AttributeTok{{-}{-}no{-}audio} \AttributeTok{{-}f} \AttributeTok{{-}w}

\CommentTok{\# Grabación profesional}
\ExtensionTok{scrcpy} \AttributeTok{{-}{-}video{-}codec}\OperatorTok{=}\NormalTok{h265 }\AttributeTok{{-}m1920} \AttributeTok{{-}{-}max{-}fps}\OperatorTok{=}\NormalTok{60 }\DataTypeTok{\textbackslash{}}
  \AttributeTok{{-}{-}audio{-}codec}\OperatorTok{=}\NormalTok{aac }\AttributeTok{{-}r}\NormalTok{ tutorial.mp4}

\CommentTok{\# Control inalámbrico ligero}
\ExtensionTok{scrcpy} \AttributeTok{{-}{-}tcpip} \AttributeTok{{-}m1024} \AttributeTok{{-}{-}max{-}fps}\OperatorTok{=}\NormalTok{30}

\CommentTok{\# Demo con toques visibles}
\ExtensionTok{scrcpy} \AttributeTok{{-}m1920} \AttributeTok{{-}{-}show{-}touches} \AttributeTok{{-}{-}stay{-}awake}

\CommentTok{\# Solo grabar micrófono}
\ExtensionTok{scrcpy} \AttributeTok{{-}{-}audio{-}source}\OperatorTok{=}\NormalTok{mic }\AttributeTok{{-}{-}no{-}video} \DataTypeTok{\textbackslash{}}
  \AttributeTok{{-}{-}no{-}playback} \AttributeTok{{-}r}\NormalTok{ grabacion.opus}
\end{Highlighting}
\end{Shaded}

\section{Configuración Permanente}\label{configuraciuxf3n-permanente}

\subsection{Alias en Fish}\label{alias-en-fish}

\textbf{Agregar a \texttt{\textasciitilde{}/.config/fish/config.fish}:}

\begin{Shaded}
\begin{Highlighting}[]
\NormalTok{\# =============================================================================}
\NormalTok{\# ALIAS DE SCRCPY}
\NormalTok{\# =============================================================================}

\NormalTok{\# Uso básico}
\NormalTok{alias scr="scrcpy"}
\NormalTok{alias scrf="scrcpy {-}f"}
\NormalTok{alias scrw="scrcpy {-}w"}

\NormalTok{\# Presentaciones}
\NormalTok{alias scr{-}present="scrcpy {-}m1920 {-}{-}max{-}fps=60 {-}{-}no{-}audio {-}f {-}w"}
\NormalTok{alias scr{-}demo="scrcpy {-}m1920 {-}{-}show{-}touches {-}{-}stay{-}awake {-}f"}

\NormalTok{\# Grabación}
\NormalTok{alias scr{-}rec="scrcpy {-}{-}video{-}codec=h265 {-}m1920 {-}{-}max{-}fps=60 {-}{-}record"}
\NormalTok{alias scr{-}rec{-}audio="scrcpy {-}{-}audio{-}source=mic {-}{-}no{-}video {-}{-}no{-}playback {-}{-}record"}

\NormalTok{\# WiFi}
\NormalTok{alias scr{-}wifi="scrcpy {-}{-}tcpip"}

\NormalTok{\# Control}
\NormalTok{alias scr{-}control="scrcpy {-}{-}no{-}video {-}{-}no{-}audio {-}KM"}
\end{Highlighting}
\end{Shaded}

\textbf{Recargar:}

\begin{Shaded}
\begin{Highlighting}[]
\NormalTok{source \textasciitilde{}/.config/fish/config.fish}
\end{Highlighting}
\end{Shaded}

\subsection{Script de Conexión
Automática}\label{script-de-conexiuxf3n-automuxe1tica}

\textbf{Crear \texttt{\textasciitilde{}/bin/scrcpy-auto}:}

\begin{Shaded}
\begin{Highlighting}[]
\CommentTok{\#!/bin/bash}

\CommentTok{\# Script para conectar scrcpy automáticamente}

\CommentTok{\# Verificar dispositivos}
\VariableTok{DEVICES}\OperatorTok{=}\VariableTok{$(}\ExtensionTok{adb}\NormalTok{ devices }\KeywordTok{|} \FunctionTok{grep} \AttributeTok{{-}w}\NormalTok{ device }\KeywordTok{|} \FunctionTok{wc} \AttributeTok{{-}l}\VariableTok{)}

\ControlFlowTok{if} \BuiltInTok{[} \StringTok{"}\VariableTok{$DEVICES}\StringTok{"} \OtherTok{{-}eq}\NormalTok{ 0 }\BuiltInTok{]}\KeywordTok{;} \ControlFlowTok{then}
    \BuiltInTok{echo} \StringTok{"No hay dispositivos conectados."}
    \BuiltInTok{echo} \StringTok{"Conecta un dispositivo por USB o configura WiFi."}
    \BuiltInTok{exit}\NormalTok{ 1}
\ControlFlowTok{elif} \BuiltInTok{[} \StringTok{"}\VariableTok{$DEVICES}\StringTok{"} \OtherTok{{-}eq}\NormalTok{ 1 }\BuiltInTok{]}\KeywordTok{;} \ControlFlowTok{then}
    \BuiltInTok{echo} \StringTok{"Conectando al único dispositivo..."}
    \ExtensionTok{scrcpy} \AttributeTok{{-}m1920} \AttributeTok{{-}{-}max{-}fps}\OperatorTok{=}\NormalTok{60 }\StringTok{"}\VariableTok{$@}\StringTok{"}
\ControlFlowTok{else}
    \BuiltInTok{echo} \StringTok{"Múltiples dispositivos detectados:"}
    \ExtensionTok{adb}\NormalTok{ devices}
    \BuiltInTok{echo} \StringTok{""}
    \BuiltInTok{echo} \StringTok{"Usa: scrcpy {-}s SERIAL"}
\ControlFlowTok{fi}
\end{Highlighting}
\end{Shaded}

\textbf{Hacer ejecutable:}

\begin{Shaded}
\begin{Highlighting}[]
\FunctionTok{chmod}\NormalTok{ +x \textasciitilde{}/bin/scrcpy{-}auto}
\end{Highlighting}
\end{Shaded}

\textbf{Usar:}

\begin{Shaded}
\begin{Highlighting}[]
\ExtensionTok{scrcpy{-}auto}
\ExtensionTok{scrcpy{-}auto} \AttributeTok{{-}{-}record}\OperatorTok{=}\NormalTok{video.mp4}
\end{Highlighting}
\end{Shaded}

\section{Recursos Adicionales}\label{recursos-adicionales}

\subsection{Documentación Oficial}\label{documentaciuxf3n-oficial}

\begin{itemize}
\tightlist
\item
  \textbf{GitHub:} https://github.com/Genymobile/scrcpy
\item
  \textbf{Documentación:}
  https://github.com/Genymobile/scrcpy/tree/master/doc
\item
  \textbf{FAQ:} https://github.com/Genymobile/scrcpy/blob/master/FAQ.md
\end{itemize}

\subsection{Comunidad}\label{comunidad}

\begin{itemize}
\tightlist
\item
  \textbf{Reddit:} r/scrcpy
\item
  \textbf{Issues:} https://github.com/Genymobile/scrcpy/issues
\end{itemize}

\subsection{Herramientas
Complementarias}\label{herramientas-complementarias}

\textbf{Para editar grabaciones:}

\begin{Shaded}
\begin{Highlighting}[]
\FunctionTok{sudo}\NormalTok{ pacman }\AttributeTok{{-}S}\NormalTok{ ffmpeg kdenlive obs{-}studio}
\end{Highlighting}
\end{Shaded}

\textbf{Para capturas de pantalla adicionales:}

\begin{Shaded}
\begin{Highlighting}[]
\FunctionTok{sudo}\NormalTok{ pacman }\AttributeTok{{-}S}\NormalTok{ flameshot}
\end{Highlighting}
\end{Shaded}

\textbf{Para streaming:}

\begin{Shaded}
\begin{Highlighting}[]
\FunctionTok{sudo}\NormalTok{ pacman }\AttributeTok{{-}S}\NormalTok{ obs{-}studio}
\end{Highlighting}
\end{Shaded}

\section{Resumen de Mejores
Prácticas}\label{resumen-de-mejores-pruxe1cticas}

\subsection{Para Presentaciones}\label{para-presentaciones}

\begin{Shaded}
\begin{Highlighting}[]
\ExtensionTok{scrcpy} \AttributeTok{{-}m1920} \AttributeTok{{-}{-}max{-}fps}\OperatorTok{=}\NormalTok{60 }\AttributeTok{{-}{-}no{-}audio} \AttributeTok{{-}{-}always{-}on{-}top} \AttributeTok{{-}{-}stay{-}awake} \AttributeTok{{-}f}
\end{Highlighting}
\end{Shaded}

\subsection{Para Grabaciones}\label{para-grabaciones}

\begin{Shaded}
\begin{Highlighting}[]
\ExtensionTok{scrcpy} \AttributeTok{{-}{-}video{-}codec}\OperatorTok{=}\NormalTok{h265 }\AttributeTok{{-}m1920} \AttributeTok{{-}{-}max{-}fps}\OperatorTok{=}\NormalTok{60 }\DataTypeTok{\textbackslash{}}
  \AttributeTok{{-}{-}audio{-}codec}\OperatorTok{=}\NormalTok{aac }\AttributeTok{{-}{-}show{-}touches} \DataTypeTok{\textbackslash{}}
  \AttributeTok{{-}{-}record}\OperatorTok{=}\NormalTok{archivo.mp4}
\end{Highlighting}
\end{Shaded}

\subsection{Para Control Remoto}\label{para-control-remoto}

\begin{Shaded}
\begin{Highlighting}[]
\CommentTok{\# Primera vez (USB)}
\ExtensionTok{scrcpy} \AttributeTok{{-}{-}tcpip}

\CommentTok{\# Siguientes veces (WiFi)}
\ExtensionTok{scrcpy} \AttributeTok{{-}m1024} \AttributeTok{{-}{-}max{-}fps}\OperatorTok{=}\NormalTok{30}
\end{Highlighting}
\end{Shaded}

\subsection{Para Documentación}\label{para-documentaciuxf3n}

\begin{Shaded}
\begin{Highlighting}[]
\ExtensionTok{scrcpy} \AttributeTok{{-}m1920} \AttributeTok{{-}{-}show{-}touches} \AttributeTok{{-}{-}stay{-}awake} \AttributeTok{{-}{-}record}\OperatorTok{=}\NormalTok{doc.mkv}
\end{Highlighting}
\end{Shaded}

\section{Publicaciones Similares}\label{publicaciones-similares}

Si te interesó este artículo, te recomendamos que explores otros blogs y
recursos relacionados que pueden ampliar tus conocimientos. Aquí te dejo
algunas sugerencias:

\begin{enumerate}
\def\labelenumi{\arabic{enumi}.}
\tightlist
\item
  \href{https://chaska-x.netlify.app/operating-system/2017-05-21-comandos-de-informacion-windows/index.pdf}{\faIcon{file-pdf}}
  \href{https://chaska-x.netlify.app/operating-system/2017-05-21-comandos-de-informacion-windows}{Comandos
  De Informacion Windows}
\item
  \href{https://chaska-x.netlify.app/operating-system/2019-06-19-adb/index.pdf}{\faIcon{file-pdf}}
  \href{https://chaska-x.netlify.app/operating-system/2019-06-19-adb}{Adb}
\item
  \href{https://chaska-x.netlify.app/operating-system/2021-08-17-limpieza-y-optimizacion-de-pc/index.pdf}{\faIcon{file-pdf}}
  \href{https://chaska-x.netlify.app/operating-system/2021-08-17-limpieza-y-optimizacion-de-pc}{Limpieza
  Y Optimizacion De Pc}
\item
  \href{https://chaska-x.netlify.app/operating-system/2021-10-21-scrcpy/index.pdf}{\faIcon{file-pdf}}
  \href{https://chaska-x.netlify.app/operating-system/2021-10-21-scrcpy}{Scrcpy}
\item
  \href{https://chaska-x.netlify.app/operating-system/2022-05-12-gestionar-versiones-de-jdk-en-kubuntu/index.pdf}{\faIcon{file-pdf}}
  \href{https://chaska-x.netlify.app/operating-system/2022-05-12-gestionar-versiones-de-jdk-en-kubuntu}{Gestionar
  Versiones De Jdk En Kubuntu}
\item
  \href{https://chaska-x.netlify.app/operating-system/2022-07-21-instalar-tor-browser/index.pdf}{\faIcon{file-pdf}}
  \href{https://chaska-x.netlify.app/operating-system/2022-07-21-instalar-tor-browser}{Instalar
  Tor Browser}
\item
  \href{https://chaska-x.netlify.app/operating-system/2022-08-14-crear-enlaces-duros-o-hard-link-en-linux/index.pdf}{\faIcon{file-pdf}}
  \href{https://chaska-x.netlify.app/operating-system/2022-08-14-crear-enlaces-duros-o-hard-link-en-linux}{Crear
  Enlaces Duros O Hard Link En Linux}
\item
  \href{https://chaska-x.netlify.app/operating-system/2022-09-27-comandos-vim/index.pdf}{\faIcon{file-pdf}}
  \href{https://chaska-x.netlify.app/operating-system/2022-09-27-comandos-vim}{Comandos
  Vim}
\item
  \href{https://chaska-x.netlify.app/operating-system/2023-02-16-guia-de-git-y-github/index.pdf}{\faIcon{file-pdf}}
  \href{https://chaska-x.netlify.app/operating-system/2023-02-16-guia-de-git-y-github}{Guia
  De Git Y Github}
\item
  \href{https://chaska-x.netlify.app/operating-system/2023-05-02-00-primeros-pasos-en-linux/index.pdf}{\faIcon{file-pdf}}
  \href{https://chaska-x.netlify.app/operating-system/2023-05-02-00-primeros-pasos-en-linux}{00
  Primeros Pasos En Linux}
\item
  \href{https://chaska-x.netlify.app/operating-system/2023-06-17-01-introduccion-linux/index.pdf}{\faIcon{file-pdf}}
  \href{https://chaska-x.netlify.app/operating-system/2023-06-17-01-introduccion-linux}{01
  Introduccion Linux}
\item
  \href{https://chaska-x.netlify.app/operating-system/2023-06-18-02-distribuciones-linux/index.pdf}{\faIcon{file-pdf}}
  \href{https://chaska-x.netlify.app/operating-system/2023-06-18-02-distribuciones-linux}{02
  Distribuciones Linux}
\item
  \href{https://chaska-x.netlify.app/operating-system/2023-06-19-03-instalacion-linux/index.pdf}{\faIcon{file-pdf}}
  \href{https://chaska-x.netlify.app/operating-system/2023-06-19-03-instalacion-linux}{03
  Instalacion Linux}
\item
  \href{https://chaska-x.netlify.app/operating-system/2023-06-20-04-administracion-particiones-volumenes/index.pdf}{\faIcon{file-pdf}}
  \href{https://chaska-x.netlify.app/operating-system/2023-06-20-04-administracion-particiones-volumenes}{04
  Administracion Particiones Volumenes}
\item
  \href{https://chaska-x.netlify.app/operating-system/2023-07-01-atajos-de-teclado-y-comandos-para-usar-vim/index.pdf}{\faIcon{file-pdf}}
  \href{https://chaska-x.netlify.app/operating-system/2023-07-01-atajos-de-teclado-y-comandos-para-usar-vim}{Atajos
  De Teclado Y Comandos Para Usar Vim}
\item
  \href{https://chaska-x.netlify.app/operating-system/2024-07-15-instalando-specitify/index.pdf}{\faIcon{file-pdf}}
  \href{https://chaska-x.netlify.app/operating-system/2024-07-15-instalando-specitify}{Instalando
  Specitify}
\item
  \href{https://chaska-x.netlify.app/operating-system/2025-07-10-gestiona-tus-dotfiles-con-gnu-stow/index.pdf}{\faIcon{file-pdf}}
  \href{https://chaska-x.netlify.app/operating-system/2025-07-10-gestiona-tus-dotfiles-con-gnu-stow}{Gestiona
  Tus Dotfiles Con Gnu Stow}
\item
  \href{https://chaska-x.netlify.app/operating-system/2025-12-30-guia-de-ranger/index.pdf}{\faIcon{file-pdf}}
  \href{https://chaska-x.netlify.app/operating-system/2025-12-30-guia-de-ranger}{Guia
  De Ranger}
\item
  \href{https://chaska-x.netlify.app/operating-system/2025-12-31-guia-de-kitty-terminal/index.pdf}{\faIcon{file-pdf}}
  \href{https://chaska-x.netlify.app/operating-system/2025-12-31-guia-de-kitty-terminal}{Guia
  De Kitty Terminal}
\item
  \href{https://chaska-x.netlify.app/operating-system/2025-12-31-guia-de-positron-ide/index.pdf}{\faIcon{file-pdf}}
  \href{https://chaska-x.netlify.app/operating-system/2025-12-31-guia-de-positron-ide}{Guia
  De Positron Ide}
\item
  \href{https://chaska-x.netlify.app/operating-system/2026-04-23-guia-de-rsync/index.pdf}{\faIcon{file-pdf}}
  \href{https://chaska-x.netlify.app/operating-system/2026-04-23-guia-de-rsync}{Guia
  De Rsync}
\end{enumerate}

Esperamos que encuentres estas publicaciones igualmente interesantes y
útiles. ¡Disfruta de la lectura!






\end{document}
